\documentclass[10pt,reqno]{amsart}

\usepackage[T1]{fontenc}
\usepackage{amsmath}
\usepackage{amssymb}
\usepackage{mathtools}
\usepackage{booktabs}
\usepackage{array}
\usepackage{enumitem}
\usepackage{hyperref}
\usepackage{url}
\def\UrlBreaks{\do\_\do\/\do\-}

\hypersetup{
  hidelinks,
  pdftitle={Auditing a Scoped Calculus of Equality Traces: Structural Metatheory, Contextual Reduction, and Derivation Erasure},
  pdfauthor={Arthur Freitas Ramos, Ruy J. G. B. de Queiroz, Anjolina Grisi de Oliveira, and Dov M. Gabbay}
}
\allowdisplaybreaks
\emergencystretch=3em
\setlength{\abovedisplayskip}{7pt plus 2pt minus 3pt}
\setlength{\belowdisplayskip}{7pt plus 2pt minus 3pt}
\setlength{\abovedisplayshortskip}{4pt plus 2pt}
\setlength{\belowdisplayshortskip}{5pt plus 2pt minus 2pt}

% Archive note: replace or supplement the stable commit citation with the
% future Zenodo archival DOI in the submission version.
\title[Auditing a Scoped Calculus of Equality Traces]
  {Auditing a Scoped Calculus of Equality Traces:\\
   Structural Metatheory, Contextual Reduction,
   and Derivation Erasure}

\author{Arthur Freitas Ramos}
\author{Ruy J. G. B. de Queiroz}
\author{Anjolina Grisi de Oliveira}
\author{Dov M. Gabbay}
\renewcommand{\shortauthors}{Ramos et al.}

\subjclass[2020]{03B15, 03B38, 03B70, 18N45}
\keywords{dependent type theory, de Bruijn syntax, definitional equality,
substitution, syntactic quotient, equality traces, path induction}

\newtheorem{theorem}{Theorem}[section]
\newtheorem{proposition}[theorem]{Proposition}
\newtheorem{lemma}[theorem]{Lemma}
\newtheorem{corollary}[theorem]{Corollary}
\newtheorem{definition}[theorem]{Definition}
\newtheorem{example}[theorem]{Example}
\newtheorem{remark}[theorem]{Remark}
\newtheorem{warning}[theorem]{Warning}

\newcommand{\Fin}{\mathsf{Fin}}
\newcommand{\Expr}{\mathsf{Expr}}
\newcommand{\Ctx}{\mathsf{Ctx}}
\newcommand{\Ext}[2]{\mathsf{ext}(#1,#2)}
\newcommand{\EmptyCtx}{\mathsf{empty}}
\newcommand{\Ren}{\mathsf{Ren}}
\newcommand{\Sub}{\mathsf{Sub}}
\newcommand{\Var}{\mathsf{var}}
\newcommand{\Sort}{\mathsf{U}}
\newcommand{\PiTy}{\mathsf{Pi}}
\newcommand{\SigmaTy}{\mathsf{Sigma}}
\newcommand{\Lam}{\mathsf{lam}}
\newcommand{\App}{\mathsf{app}}
\newcommand{\Pair}{\mathsf{pair}}
\newcommand{\Fst}{\mathsf{fst}}
\newcommand{\Snd}{\mathsf{snd}}
\newcommand{\NatTy}{\mathsf{Nat}}
\newcommand{\Zero}{\mathsf{zero}}
\newcommand{\Succ}{\mathsf{succ}}
\newcommand{\NatElim}{\mathsf{natElim}}
\newcommand{\CodeNat}{\widehat{\mathsf{Nat}}}
\newcommand{\CodePi}{\widehat{\mathsf{Pi}}}
\newcommand{\CodeSigma}{\widehat{\mathsf{Sigma}}}
\newcommand{\CodeId}{\widehat{\mathsf{Id}}}
\newcommand{\El}{\mathsf{El}}
\newcommand{\IdTy}{\mathsf{Id}}
\newcommand{\Refl}{\mathsf{refl}}
\newcommand{\EqJ}{\mathsf{eqJ}}
\newcommand{\Tr}{\mathsf{Tr}}
\newcommand{\TrStep}{\mathrel{\triangleright_{\mathrm{tr}}}}
\newcommand{\TrEq}{\mathrel{\simeq_{\mathrm{tr}}}}
\newcommand{\Frame}{\mathsf{Frame}}
\newcommand{\IProg}{\mathsf{IProg}}
\newcommand{\ev}{\mathsf{ev}}
\newcommand{\den}[1]{\mathopen{\lbrack\!\lbrack}#1\mathclose{\rbrack\!\rbrack}}
\newcommand{\defeq}{\mathrel{\equiv_{\mathrm d}}}
\newcommand{\red}{\mathrel{\longmapsto}}
\newcommand{\ctxred}{\mathrel{\rightarrowtail}}
\newcommand{\vdashC}{\vdash_{\mathrm{c}}}
\newcommand{\QProg}{\mathsf{QProg}}
\newcommand{\Framed}{\mathsf{Framed}}
\newcommand{\ctxredstar}{\mathrel{\rightarrowtail^{*}}}
\newcommand{\srw}{\mathrel{\rightsquigarrow_{\mathrm s}}}
\newcommand{\srweq}{\mathrel{\simeq_{\mathrm s}}}
\newcommand{\idren}{\mathsf{id}}
\newcommand{\wk}{\mathsf{wk}}
\newcommand{\lift}{\uparrow}
\newcommand{\Q}{\mathcal{Q}}

\begin{document}

\begin{abstract}
This is a machine-checked audit, not a new type theory.  Its subject is the raw
scope-indexed calculus that underlies a formalized development of equality
traces, and its output is a precise account of what that calculus proves, what
it fails to prove, and what its traces actually record.  We report the failures
as theorems rather than as gaps, which is what distinguishes an audit from an
incomplete metatheory.

The calculus is a raw, scope-indexed dependent-style syntax in de Bruijn form,
with dependent products and sums, level-indexed natural numbers, Tarski-style
codes, identity, reflexivity, and a two-binder raw \(J\)-constructor.  Contexts
are same-scope assignments rather than well-formed telescopes; definitional
equality \(\defeq\) is untyped; and there is no conversion rule.  Renaming and
simultaneous substitution are given explicitly under one and two binders, with
their identity, composition, lifting, weakening, and instantiation laws, and
substitution is proved admissible for the raw typing judgment.  Ten top-level
primitive redexes preserve typing modulo \(\defeq\) between the displayed type
expressions, and the congruence closure of those rules, together with its
reflexive-transitive closure, is sound for \(\defeq\) and stable under
substitution.

Contextual subject reduction then \emph{fails}, and we prove that it does: a
reduction inside a well-typed term yields a reduct admitting no derivation at
the same type.  We prove further that neither rule usually invoked to repair
this---conversion and its context-wise companion---is admissible, so the natural
hypothesis record packaging them is \emph{uninhabited} and any conditional
subject-reduction theorem over it is vacuous as a statement about the raw
judgment.  What such a proof does establish is an analysis of the induction:
repairs of exactly those two shapes close every case.  The conversion-closed
extension satisfies both rules by construction, and we state exactly what an
unconditional theorem for it would additionally need---generation lemmas, hence
confluence---which we do not prove.

The audit's second half concerns what the traces retain.  Substitution descends
to a representative-independent, compositional action on the syntactic quotient
\(\Q_n=\Expr_n/{\defeq}\), and an independent source language of identity
programs evaluates into \emph{unlabelled quotient traces}: lists of equality
steps over \(\Q_n\) together with an endpoint equality.  That these traces are
unlabelled is proved, not asserted.  Evaluation \emph{factors} through a
label-free syntax whose atoms carry only proof-irrelevant equalities of classes
and which therefore cannot mention a source derivation; program size and
constructor shape are preserved, so the arrangement of steps survives while the
reasons do not; and every quotient equality between represented classes has a
source witness erasing back to the same atom.  Erasure is \emph{not} surjective
onto the label-free grammar, whose congruence constructor accepts arbitrary
endofunctions, so we characterize its essential image exactly---the fragment
whose atoms are quotient soundness applied to a source \(\defeq\)-derivation and
whose congruences are induced by source frames.  This separates the present
target sharply from label-sensitive presentations of computational paths.

The verdict is deflationary and quantitative.  A trace is observable list
metadata decorating an equality: its entries need not be adjacent, composable,
or related to the outer endpoints, and the erasure theorem says the reasons have
been discarded.  Applying the metadata-fiber criterion of the companion article
\cite{ObservableMetadata2026} to our own records computes the trace fiber to be
a list of quotient classes, which is never contractible, so trace records admit
no unrestricted based eliminator, while motives factoring through the equality
projection remain available.  No judgmental beta is claimed, and no model of
Martin-L\"of or homotopy type theory is supplied.  A Lean~4 development
machine-checks every statement above: the structural laws, the reduction theory,
the failure and inadmissibility theorems, quotient substitution, erasure and its
essential image, rewrite soundness, and the trace obstruction.
\end{abstract}

\maketitle

\section{Introduction}
\label{sec:introduction}

This paper audits a formalized calculus of equality traces.  It is not a
proposal for a new type theory, and it is not an interpretation theorem.  The
object under examination is a raw, scope-indexed dependent-style syntax that
exists as Lean~4 source, and the task is to determine, with machine-checked
precision, three things: which structural properties it genuinely has, which
properties it demonstrably lacks, and what its equality traces actually record.
Where the calculus fails, we prove the failure.  That discipline---negative
results stated as theorems rather than reported as gaps---is what an audit
contributes, and it is what a reader deciding whether to build on the artifact
needs.

The framing matters because raw calculi are easy to over-read.  A syntax with
\(\Pi\)-types, universes, an identity type, and a \(J\)-constructor invites the
assumption that the usual metatheory is present in the background.  Here it is
not, and the paper's job is to say exactly where the resemblance stops.

\subsection*{The object of study}

Fix \(n\in\mathbb N\).  The
expressions \(\Expr_n\) have exactly \(n\) de Bruijn variables and the
constructors listed in Section~\ref{sec:syntax}.  A context is merely a
same-scope assignment
\(\Gamma:\Fin(n)\to\Expr_n\).  There is no context-formation judgment saying
that \(\Gamma\) is a well-formed telescope.  The typing family
\(\Gamma\vdash t:A\) is raw; its definitional equality
\(t\defeq u\) is an untyped relation on expressions.  There is no conversion
rule and no typed definitional-equality judgment.

Let \(\Q_n=\Expr_n/{\defeq}\) be the syntactic quotient and write
\(\den t\) for the class of \(t\).  Over any carrier \(X\) we use the record
\[
 \Tr_X(a,b)
   =\operatorname{List}\!\left(\sum_{x,y:X}(x=y)\right)\times(a=b).
\tag{1.1}\label{eq:intro-trace-record}
\]
When \(X=\Q_n\), this is an \emph{unlabelled quotient trace}.  Its list
contains equality steps over already-quotiented endpoints.  In particular,
a source derivation of \(t\defeq u\) is not retained as a label: it is sent to
the proof-irrelevant quotient equality \(\den t=\den u\).
Throughout, \(X\simeq Y\) denotes specified mutually inverse maps, and
\(\mathsf{transport}_D(e,-)\) denotes equality transport in a family \(D\);
Section~\ref{sec:ambient} states the precise universe conventions and
propositional computation law.

Two questions organize the audit.  The first is how far the raw calculus
gets on its own, and where exactly it stops.  The second is what its traces
actually record.

\subsection*{Reduction, and what conversion actually costs}

The primitive computation rules preserve typing.  That is a statement about
\emph{top-level} redexes only, and by itself it says nothing about reducing
inside a term or reducing repeatedly.  We close both gaps.
Section~\ref{sec:reduction} defines the congruence closure \(\ctxred\) of the
primitive rules, enumerated by raw constructor so the supported closure stays
auditable, together with its reflexive-transitive closure \(\ctxredstar\); both
are sound for \(\defeq\) and stable under simultaneous substitution.

Contextual subject reduction then fails, and we prove that it fails rather than
reporting that a proof gets stuck.

\begin{theorem}[Failure of contextual subject reduction]
\label{thm:intro-sr-fails}
There are a context \(\Gamma\), expressions \(t\ctxred u\), and a type \(A\)
with \(\Gamma\vdash t:A\) derivable and \(\Gamma\vdash u:A\) underivable.
\end{theorem}

The witness is small enough to check by eye: apply a variable whose
\(\Pi\)-type has the bound variable itself as codomain, so the type of the
application \emph{is} its argument, and reduce the argument.  Because raw contexts
are arbitrary same-scope assignments, no context well-formedness obligation is
incurred, and because the raw typing rules are syntax-directed, the negative
half is inversion.

The usual repair is to add a conversion rule retyping along \(\defeq\), together
with a context-conversion rule retyping along pointwise \(\defeq\) for the case
where a domain annotation is reduced under a binder.  Both are genuinely absent
here, and again this is a theorem.

\begin{theorem}[Neither rule is admissible]
\label{thm:intro-not-admissible}
The raw judgment does not admit conversion: \(\mathsf{zero}\) is typed at
\(\mathsf{Nat}\) and at nothing definitionally equal to it.  It does not admit
context conversion either: there is a family that is a type over a domain and
not over a definitionally equal one.  Consequently the record packaging the two
rules is uninhabited.
\end{theorem}

Theorem~\ref{thm:intro-not-admissible} has a consequence that is easy to state
misleadingly, so we state it plainly.  Assuming the two rules as a hypothesis
record and deriving contextual and multi-step subject reduction from it---which
we do, in Section~\ref{sec:reduction}---yields a theorem with an empty
hypothesis.

\begin{theorem}[Conditional subject reduction, and what it means]
\label{thm:intro-subject-reduction}
Assume the two omitted structural rules as an explicit hypothesis record.  Then
for every \(t\ctxred u\) and every \(\Gamma\vdash t:A\) one has
\(\Gamma\vdash u:A\), at the same type, for the full congruence closure; and the
same holds for \(t\ctxredstar u\).  By
Theorem~\ref{thm:intro-not-admissible} the hypothesis is unsatisfiable for
\(\vdash\), so the theorem is vacuous as a statement \emph{about} \(\vdash\).
Its content is the analysis of the induction: repairs of exactly those two
shapes, and no others, are needed to close every one of its cases.
\end{theorem}

Accordingly we do not say that the two rules are ``all that is missing''
between this calculus and a conversion-closed one; that phrasing confuses a fact
about a proof with a fact about a judgment.  The defensible statement is that
conversion and context conversion are \emph{sufficient to repair every
obstruction arising in the structural subject-reduction proof}, and
\emph{necessary} in the concrete sense of
Theorem~\ref{thm:intro-not-admissible}.

The positive route is to change the judgment rather than to assume rules about
it.  Section~\ref{sec:reduction} defines the conversion-closed extension
\(\vdashC\), in which both rules hold by construction and which types the reduct
that defeated \(\vdash\) in Theorem~\ref{thm:intro-sr-fails}.  We do not prove
unconditional contextual subject reduction for \(\vdashC\), and we are explicit
about why: inverting a conversion-closed judgment needs generation lemmas, and
generation lemmas for an untyped definitional equality need confluence, which
this development does not establish.  The distance from the raw calculus to
subject reduction is thus not ``two hypotheses'' but ``a different judgment,
whose metatheory needs confluence''.

\subsection*{Traces are unlabelled, and this is a theorem}

Let \(\Q_n\) be the syntactic quotient.  A source derivation of \(t\defeq u\)
is used only to obtain the equality \(\den t=\den u\), which is
proof-irrelevant; which derivation was supplied is not retained.  Since the
computational-path literature conventionally does the opposite---equality
reasons are first-class labels---this difference deserves proof rather than
assertion.

Section~\ref{sec:erasure} supplies one.  We introduce a label-free program
syntax over the quotient whose atoms carry only an equality of classes, and
which therefore has no room to mention a source derivation.

\begin{theorem}[Erasure factorization]
\label{thm:intro-erasure}
Evaluation of source identity programs factors through erasure:
\[
 \ev(p)=\ev_{\Q}(\lfloor p\rfloor)
 \qquad\text{for every source program }p,
\tag{1.2}\label{eq:intro-erasure}
\]
where \(\lfloor-\rfloor\) discards every source derivation.  Moreover
\(\lfloor-\rfloor\) preserves program size, and at atoms it is inverse to the
recovery map supplied by quotient exactness.
\end{theorem}

The three clauses say three different things.  Factorization says the
label-free syntax is sufficient: no target trace can depend on which derivation
an atom supplied.  Size preservation says erasure is not a collapse---the
arrangement of equality steps survives, and only labels are discarded.  The
round trip is \emph{atom-level}: every quotient equality between represented
classes has a source witness, and erasing that witness returns the same atom.

The three clauses do not say that the two program languages are equivalent.  The
label-free grammar is strictly more permissive than erasure needs: its
congruence constructor accepts an arbitrary endofunction of the term model,
whereas erasure only ever supplies one induced by a source frame.
Section~\ref{sec:erasure} therefore characterizes the essential image instead.

\begin{theorem}[Essential image of erasure]
\label{thm:intro-image}
Call a label-free program \emph{framed} when its atoms are quotient soundness
applied to a source definitional equality and its congruences are induced by
source frames.  Then a label-free program between represented classes is an
erasure if and only if it is framed.
\end{theorem}

Theorem~\ref{thm:intro-image} is what makes the factorization a statement about
a genuine target language rather than about an ambient grammar that happens to
contain the image.

\begin{theorem}[Raw scoped structural results]
\label{thm:intro-main}
For the raw scoped calculus with \(\Pi\), \(\Sigma\), level-indexed natural
syntax, Tarski-style codes, identity, reflexivity, and raw two-binder
\(\EqJ\):
\begin{enumerate}
\item renamings and simultaneous substitutions have explicit one- and
  two-binder lifts satisfying identity, composition, weakening, and one- and
  two-variable instantiation;
\item pointwise typed substitution preserves the raw typing family, and each of
  the ten top-level primitive computation rules preserves typing modulo
  \(\defeq\) between the displayed type expressions;
\item \(\ctxred\) and \(\ctxredstar\) are sound for \(\defeq\) and stable under
  substitution, while contextual subject reduction fails
  (Theorem~\ref{thm:intro-sr-fails});
\item \(\defeq\) is stable under substitution and under pointwise
  \(\defeq\)-equal environments, so every \(\sigma:\Sub(n,m)\) induces a
  representative-independent \(\Q(\sigma):\Q_n\to\Q_m\) respecting identity and
  composition;
\item the source family \(\IProg_n(t,u)\) is endpoint-trivial, has the expected
  unique structural-recursion interpretation, and its named rewrite equivalence
  is sound for the target one, \(p\srweq q\Rightarrow\ev(p)\TrEq\ev(q)\);
\item evaluation factors through a label-free quotient syntax whose essential
  image is characterized by Theorem~\ref{thm:intro-image}; and
\item raw \(\EqJ\) has its listed syntactic reflexivity beta constructor, not
  obtained from and not packaged with any external interpretation of its
  motive.
\end{enumerate}
\end{theorem}

\begin{proposition}[Endpoint exactness of the syntactic quotient]
\label{prop:intro-endpoint}
For \(t,u\in\Expr_n\),
\[
 \operatorname{Nonempty}\bigl(\Tr_{\Q_n}(\den t,\den u)\bigr)
 \quad\Longleftrightarrow\quad t\defeq u.
\tag{1.3}\label{eq:intro-endpoint}
\]
\end{proposition}

This equivalence is tautological after quotient construction.  The forward
direction forgets the list and uses quotient exactness; the reverse direction
uses quotient soundness and forms a singleton trace.  It is included to
identify exactly what endpoint statement is available, not as a semantic
completeness theorem.

\begin{corollary}[Trace-metadata instance]
\label{thm:intro-obstruction}
The record \(\Tr_X\) admits an eliminator through its ambient-equality
component.  It does not admit an unrestricted based eliminator satisfying
propositional beta for arbitrary motives that can observe the list field.
\end{corollary}

The obstruction uses a well-formed singleton reflexive trace, not a malformed
or non-composable list.  Hence imposing intrinsic composability on traces
does not change the classification.  The metadata fiber must be made
contractible, or motives must be restricted through the equality projection.

\subsection*{Contribution}

The contribution is an audit with two substantive findings and a supporting
body of verified structural metatheory.

The first finding is negative and concerns reduction.  It replaces a reported
obstruction by theorems: contextual subject reduction is false
(Theorem~\ref{thm:intro-sr-fails}), neither repairing rule is admissible, and
the hypothesis record packaging them is uninhabited
(Theorem~\ref{thm:intro-not-admissible}).  A gap in a proof and a false
statement are different objects, and only the second tells a prospective user of
the artifact that no amount of further effort will close the case.  Reporting a
conditional theorem as if it were a measurement would have concealed exactly
that.

The second finding concerns labels, and it converts a naming convention into
mathematics.  Calling a trace ``unlabelled'' is otherwise a claim about intent;
the factorization, size-preservation, and atom round-trip clauses of
Theorem~\ref{thm:intro-erasure}, together with the image characterization of
Theorem~\ref{thm:intro-image}, instead pin down what the target retains, what it
cannot express, and which of its programs are actually reached.  The
consequence for the wider programme is stated in Section~\ref{sec:comparison}:
the traces produced here are not label-sensitive computational paths, and the
difference is provable rather than stylistic.

Supporting these are binder-correct raw scoped syntax; source-derived
preservation for ten primitive redexes; a contextual and multi-step reduction
theory; representative-independent and compositional quotient substitution;
source/target-separated identity programs; and the trace-record elimination
obstruction.  Handling the two-binder constructors correctly is the least
glamorous and most error-prone part of this, and mechanizing it is what makes
the negative results trustworthy: a counterexample to subject reduction is only
as good as the substitution algebra it is stated over.

Several components are deliberately modest and we mark them as such.  The
metadata-fiber criterion is imported from the companion article
\cite{ObservableMetadata2026}; what is done here is its application to this
paper's own trace records.  Source identity rewrite soundness is a homomorphic
simulation, largely forced by the structural-recursion interpretation once the
target operations and their laws are supplied.  Endpoint exactness is
tautological after quotienting.  The erasure factorization is likewise forced
in part by the definitions---source atoms are evaluated only through quotient
soundness---so its interest lies in the image characterization rather than in
the induction.  The paper claims no absolute priority.

\subsection*{What this paper is not}

The terminology in Theorem~\ref{thm:intro-main} is intentional.  This is a
\emph{scoped dependent-style calculus}, not a presentation of full
Martin-L\"of type theory: we construct a syntactic quotient and no typed model
or initial CwF, source identity programs have no typing judgment and are
endpoint-trivial, and there is no normalization, confluence, canonicity,
decidable type checking, univalence, higher inductive type, or completeness
theorem for source identity rewrites.  Section~\ref{sec:limitations} states
these boundaries in full, and every one of them is load-bearing somewhere
above.

We also do not repair the calculus.  The conversion-closed extension of
Section~\ref{subsec:conversion-closed} shows what a repair would look like and
identifies precisely what it would then owe---generation lemmas, hence
confluence---but supplying that positive metatheory is separate work, and
presenting it here would obscure the audit's actual result, which is that the
calculus as formalized does not have it.

\subsection*{Relation to earlier computational paths}

The prior computational-path articles
\cite{deQueirozEtAl2016,RamosEtAl2017,RamosEtAl2018} and the 2026 monograph
\cite{RamosEtAl2026} study label-sensitive paths: equality reasons are
first-class labels and their rewrite calculus preserves that intensional
presentation.  The target here is weaker by design.  A \(\defeq\)-derivation is
used only to obtain equality in \(\Q_n\), and proof irrelevance then erases
which derivation was chosen, so evaluated identity programs decorate
quotient-collapsed endpoints rather than carrying the original reason.  This
change of target is central, and Section~\ref{sec:comparison} compares the two
statement by statement.

\subsection*{Organization}

Section~\ref{sec:layers} separates the equality layers.
Sections~\ref{sec:syntax} and~\ref{sec:substitution} give the raw syntax and
its structural algebra.  Section~\ref{sec:judgments} presents typing,
computation, substitution, and primitive-redex preservation.
Section~\ref{sec:reduction} develops contextual and multi-step reduction and
proves the conditional subject-reduction theorem.  The syntactic
quotient is constructed in Section~\ref{sec:syntactic-quotient}.  Source identity programs
and their trace semantics occupy Section~\ref{sec:identity-programs}, and
Section~\ref{sec:erasure} proves that those traces are unlabelled.
Section~\ref{sec:example} gives lambda, dependent-projection, and
natural-successor examples.  Section~\ref{sec:j-boundary} says what the
resulting traces are, and Section~\ref{sec:comparison} compares them with
prior work.
Appendices~\ref{app:grammar} and~\ref{app:rewrites} collect the raw rules.

\section{The equality layers}
\label{sec:layers}

\subsection{Ambient metatheory}
\label{sec:ambient}

Fix once and for all a Grothendieck-style universe \(\mathcal U\) of small
sets.  Every carrier, expression family, metadata fiber, and quotient used
below is an element of \(\mathcal U\).  The universe is closed under finite
lists, dependent sums and products indexed by small sets, equality
propositions, and quotients by equivalence relations.  We reason in ordinary
extensional mathematics about these sets.  Concretely:
\begin{enumerate}
\item equality has its usual eliminator and equality propositions are proof
  irrelevant:
  \(h,k:x=y\) implies \(h=k\);
\item functions are extensional, so pointwise equal functions are equal;
\item finite lists with append, reversal, map, and no-confusion
  \([]\neq z::L\);
\item every equivalence relation \(E\) on a small set \(X\) has an effective
  small quotient
  \(\eta:X\to X/E\) with
  \[
       E(x,y)\Longrightarrow\eta x=\eta y
       \quad\text{and}\quad
       (\eta x=\eta y)\Longrightarrow E(x,y).
  \tag{2.1}\label{eq:quotient-sound-exact}
  \]
\end{enumerate}
The first implication in~\eqref{eq:quotient-sound-exact} is quotient
soundness; the second is quotient exactness.

For small sets \(X,Y\), the notation \(X\simeq Y\) means specified maps
\[
 F:X\to Y,\qquad G:Y\to X
\]
together with propositional equalities
\(F(Gy)=y\) and \(G(Fx)=x\).  If \(D:X\to\mathcal U\) and \(e:x=y\), then
\(\mathsf{transport}_D(e,-):D(x)\to D(y)\) denotes equality transport.
We use only its propositional reflexivity law
\[
 \mathsf{transport}_D(\mathsf{refl}_x,z)=z.
\tag{2.2}\label{eq:transport-beta}
\]
No judgmental computation rule is part of the ambient assumptions.

When we quantify over the collection of all \(\mathcal U\)-valued motives
\(C\), that collection is regarded as living in one fixed larger universe
\(\mathcal V\).  Each individual motive value, equality proposition,
\(\IProg\)-fiber, \(R\)-fiber, and metadata fiber remains
\(\mathcal U\)-small.  Thus the phrase ``admits an unrestricted based
eliminator satisfying propositional beta'' means an eliminator into every
\(\mathcal U\)-valued motive together with that beta family; it makes no
impredicative claim that
the collection of all such motives lies in \(\mathcal U\) itself.
For \(X\in\mathcal U\) and \(c\in X\), the sets
\(\mathsf{Contr}_c(X)\) and \(\mathsf{isContr}(X)\) are
\(\mathcal U\)-small.  By contrast, the set of all
\(\mathcal U\)-valued motives on \(X\), the data of an eliminator over all of
them, and the proposition asserting existence of such data lie in
\(\mathcal V\).  Each individual propositional beta equation lies in the
small equality proposition of a fiber \(C(x)\in\mathcal U\).

These assumptions have sharply delimited uses.  Function extensionality
turns the pointwise lift and environment calculations of
Sections~\ref{sec:substitution} and~\ref{sec:syntactic-quotient} into
equalities of functions.  Quotient soundness defines quotient substitution
and the trace attached to a source atom.  Quotient exactness is used only in
the map from trace inhabitation back to source \(\defeq\) in endpoint
exactness.  Proof irrelevance erases the
choice of source \(\defeq\)-derivation after quotient soundness and makes the
external equality-component eliminator insensitive to equality proofs.
List no-confusion proves that the empty and singleton traces differ.
Lean is one realization of this metatheory for reproducibility; it is not the
definition of the mathematics.

\subsection{The dependency layers}

The word ``equality'' occurs at several levels.  The direction of dependence
matters: none of the target trace relations is used to generate a source
typing or definitional-equality derivation.

\begin{table}[ht]
\centering
\small
\renewcommand{\arraystretch}{1.22}
\begin{tabular}{>{\raggedright\arraybackslash}p{0.20\textwidth}
                >{\centering\arraybackslash}p{0.19\textwidth}
                >{\raggedright\arraybackslash}p{0.48\textwidth}}
\toprule
Layer & Notation & Information and role\\
\midrule
Metatheoretic equality
  & \(x=y\)
  & Proof-irrelevant equality used to form quotient classes and the endpoint
    field of a trace record.\\
Primitive source computation
  & \(t\red u\)
  & Exactly one of ten top-level constructors, with no closure of its own.\\
Contextual reduction
  & \(t\ctxred u\), \(t\ctxredstar{}u\)
  & Congruence closure of \(\red\), and its reflexive-transitive closure
    (Section~\ref{sec:reduction}).\\
Source definitional equality
  & \(t\defeq u\)
  & Equivalence and enumerated compatible closure of primitive computation.
    It is untyped and defines the syntactic quotient.\\
Syntactic-quotient equality
  & \(\den t=\den u\)
  & Ambient equality in \(\Q_n=\Expr_n/{\defeq}\).  Quotient exactness
    identifies this precisely with \(t\defeq u\).\\
Source identity program
  & \(p:\IProg_n(t,u)\)
  & Explicit syntax generated by source atoms, unit, inverse, composition,
    and finite same-scope frames.\\
Unlabelled quotient trace
  & \(T:\Tr_{\Q_n}(\den t,\den u)\)
  & A list of quotient-equality steps plus an endpoint equality.  Source
    \(\defeq\)-derivation labels have been erased.\\
Source identity rewrite
  & \(p\srweq q\)
  & Symmetric reflexive-transitive closure of the named source program laws.\\
Target trace rewrite
  & \(T\TrEq U\)
  & Explicit symmetric-reflexive-transitive trace equivalence.  Soundness maps the
    preceding row to this one; no reflection theorem is asserted.\\
\bottomrule
\end{tabular}
\caption{The layered equality discipline.}
\label{tab:equality-layers}
\end{table}

There are two essential arrows.  The first is quotient soundness
\[
       (t\defeq u)\longmapsto(\den t=\den u).
\]
The second evaluates source identity syntax:
\[
       \ev:\IProg_n(t,u)\longrightarrow
       \Tr_{\Q_n}(\den t,\den u).
\]
The first arrow erases source evidence: quotient soundness returns only an
ambient equality proof, and proof irrelevance identifies different proofs of
the same quotient equality.  The second arrow consequently cannot preserve a
\(\defeq\)-derivation as a label.

\begin{remark}[Scope of the no-target-leakage statement]
The source constructors in Definition~\ref{def:identity-programs} consume
only source \(\defeq\)-derivations and source frames.  The statement applies
to that family alone (named \texttt{RawMLTT.IdentityExpr} in the
accompanying software).  The broader software repository also contains
convenience syntax that starts from already constructed target trace records; that
syntax is excluded from the calculus.
\end{remark}

\section{An independently scoped source calculus}
\label{sec:syntax}

\subsection{Scopes and raw expressions}

Following de Bruijn's nameless representation
\cite{deBruijn1972}, let
\(\Fin(n)=\{0,\ldots,n-1\}\), with \(0\) the most recently bound
variable.  The family \(\Expr_n\) is defined inductively by
\begin{align*}
t,u,A,B,M ::= {}&
  \Var(i) \mid \Sort_\ell
  \mid \PiTy(A,B) \mid \SigmaTy(A,B)
  \mid \Lam(t) \mid \App(t,u)\\
 &\mid \Pair(t,u)\mid\Fst(t)\mid\Snd(t)
  \mid \NatTy_\ell\mid\Zero_\ell\mid\Succ_\ell(t)\\
 &\mid \NatElim_\ell(M,z,s,t)
  \mid \CodeNat_\ell\mid\CodePi(A,B)
  \mid\CodeSigma(A,B)\mid\CodeId(A,t,u)\\
 &\mid \El(A)\mid\IdTy(A,t,u)\mid\Refl(t)
  \mid \EqJ(M,r,b,p).
\end{align*}
Here \(i\in\Fin(n)\) and \(\ell\in\mathbb{N}\).  The displayed grammar
suppresses scope indices, which are part of the inductive definition:
\[
\begin{array}{c|c}
\text{constructor} & \text{scopes of arguments}\\ \hline
\PiTy,\SigmaTy,\CodePi,\CodeSigma
  & \Expr_n\times\Expr_{n+1}\to\Expr_n\\
\Lam & \Expr_{n+1}\to\Expr_n\\
\NatElim_\ell
  & \Expr_{n+1}\times\Expr_n\times\Expr_{n+2}
      \times\Expr_n\to\Expr_n\\
\EqJ
  & \Expr_{n+2}\times\Expr_n\times\Expr_n\times\Expr_n
      \to\Expr_n.
\end{array}
\]
All other arguments remain in scope \(n\).  Thus the codomains of
\(\PiTy\), \(\SigmaTy\), and their codes bind one variable; the successor
branch of \(\NatElim\) binds a predecessor and an induction hypothesis; and
the motive of \(\EqJ\) binds an endpoint and an identity proof.

A raw context of length \(n\) is a function
\[
                       \Gamma:\Fin(n)\longrightarrow\Expr_n.
\]
This is a \emph{same-scope assignment}: every declaration is already written
in the full scope \(n\).  There is no judgment
\(\vdash\Gamma\;\mathsf{ctx}\), no inductively generated telescope, and no
theorem relating these assignments to well-formed MLTT contexts.  The raw
typing constructors use whichever declarations they require.  The empty
assignment is
\(\EmptyCtx_0:\Fin(0)\to\Expr_0\).  Extension is the explicit operation
\[
 \Ext{\Gamma}{A}\in\Ctx(n+1)
 \quad(\Gamma\in\Ctx(n),\ A\in\Expr_n)
\]
that weakens every declaration:
\[
\begin{split}
\Ext{\Gamma}{A}(0)&=A[\wk_n]_{\mathrm r},\\
\Ext{\Gamma}{A}(i+1)&=\Gamma(i)[\wk_n]_{\mathrm r}.
\end{split}
\tag{3.1}\label{eq:context-extension}
\]

\subsection{Levels, natural numbers, and codes}

Levels are natural numbers internal to the grammar only as parameters.
The rules use
\[
              \Gamma\vdash\Sort_\ell:\Sort_{\ell+1},
   \qquad
              \Gamma\vdash\NatTy_\ell:\Sort_\ell.
\]
Natural numbers, zero, successor, and their recursor all carry the same
level parameter.  This indexing is not cosmetic: it makes
\[
                 \El(\CodeNat_\ell)\red\NatTy_\ell
\]
type preserving at every level.  There is no cumulativity or resizing rule.

The code constructors are Tarski-style in the limited sense relevant here:
codes are explicit syntax and are decoded by \(\El\).  A code at level
\(\ell\) has type \(\Sort_\ell\), and its decoding also has type
\(\Sort_\ell\).  Product and sum codes bind over the decoding of their domain
code.  The calculus does not purport to give a full hierarchy of inductive
universes or a universe-polymorphic elaboration algorithm.

\subsection{\texorpdfstring{The raw two-binder \(J\)-constructor}
{The raw two-binder J-constructor}}

Suppose \(A,a\in\Expr_n\).  Define the proof declaration in endpoint scope by
\[
 \mathsf{IdProof}_n(A,a)
   \coloneqq
   \IdTy\bigl(A[\wk_n]_{\mathrm r},
              a[\wk_n]_{\mathrm r},\Var(0)\bigr)
   \in\Expr_{n+1}.
\tag{3.2}\label{eq:idproof-definition}
\]
The endpoint is the new variable of \(\Ext{\Gamma}{A}\).  Extending this
assignment once more gives the explicit motive assignment
\[
 \Ext{\Ext{\Gamma}{A}}{\mathsf{IdProof}_n(A,a)}
   \in\Ctx(n+2).
\tag{3.3}\label{eq:motive-context}
\]
Because the latest variable is numbered \(0\), a motive
\(M\in\Expr_{n+2}\) sees the proof at index \(0\), the endpoint at index
\(1\), and the ambient variables from index \(2\) onward.  The raw term
\(\EqJ(M,r,b,p)\) instantiates those two variables by \(p\) and \(b\).
Its only computation constructor is the syntactic rule
\[
                 \EqJ(M,r,a,\Refl(a))\red r.
\]
Nothing in the raw syntax says that \(M\) factors through an ambient equality
component, and this paper proves no such interpretation.  The independent
external eliminator of Section~\ref{sec:j-boundary} is therefore a separate
construction, not a semantics for this raw motive.

\section{Renaming and simultaneous substitution}
\label{sec:substitution}

\subsection{Renamings and binder lift}

Set
\[
                  \Ren(n,m)\coloneqq\Fin(n)\to\Fin(m).
\]
We write \(\rho:n\rightsquigarrow_{\mathrm r}m\) only when
\(\rho\in\Ren(n,m)\).  Identity and composition are
\[
  \idren^{\Ren}_n(i)=i,\qquad
  (\rho;_{\Ren}\tau)(i)=\tau(\rho(i)),
\]
where \(\rho\in\Ren(n,m)\) and \(\tau\in\Ren(m,k)\).
Weakening \(\wk_n\in\Ren(n,n+1)\) is \(\wk_n(i)=i+1\).  The binder lift
\(\rho^\lift\in\Ren(n+1,m+1)\) is
\[
        \rho^\lift(0)=0,\qquad
        \rho^\lift(i+1)=\rho(i)+1.
\tag{4.1}\label{eq:ren-lift}
\]
Renaming \(t[\rho]_{\mathrm r}\) is structural recursion, using one lift
under one binder and two lifts under the \(\NatElim\) successor branch and
the \(\EqJ\) motive.

\begin{proposition}[Renaming algebra]
\label{prop:renaming}
For every expression and every composable pair of renamings,
\[
\begin{aligned}
t[\idren^{\Ren}_n]_{\mathrm r}&=t,\\
t[\rho]_{\mathrm r}[\tau]_{\mathrm r}
  &=t[\rho;_{\Ren}\tau]_{\mathrm r},\\
(\idren^{\Ren}_n)^\lift&=\idren^{\Ren}_{n+1},\\
(\rho;_{\Ren}\tau)^\lift&=\rho^\lift;_{\Ren}\tau^\lift,\\
\wk_n;_{\Ren}\rho^\lift&=\rho;_{\Ren}\wk_m.
\end{aligned}
\]
\end{proposition}

\begin{proof}
Prove the three equations between renamings first.  At index \(0\), both
\((\idren^{\Ren}_n)^\lift\) and
\((\rho;_{\Ren}\tau)^\lift\) return \(0\).  At \(i+1\),
\[
 (\idren^{\Ren}_n)^\lift(i+1)=i+1,
\qquad
 (\rho;_{\Ren}\tau)^\lift(i+1)
   =\tau(\rho(i))+1
   =(\rho^\lift;_{\Ren}\tau^\lift)(i+1).
\]
For \(i:\Fin(n)\),
\[
 (\wk_n;_{\Ren}\rho^\lift)(i)=\rho(i)+1
   =(\rho;_{\Ren}\wk_m)(i).
\]
Function extensionality gives the three renaming equalities.

Now induct structurally on \(t\), proving identity and composition
simultaneously.  The variable case is the corresponding equation for
\(\rho\); \(\Sort_\ell,\NatTy_\ell,\Zero_\ell,\CodeNat_\ell\) are fixed.
For a one-binder constructor, for example,
\[
\begin{aligned}
\PiTy(A,B)[\rho]_{\mathrm r}[\tau]_{\mathrm r}
 &=\PiTy(A[\rho]_{\mathrm r}[\tau]_{\mathrm r},
          B[\rho^\lift]_{\mathrm r}[\tau^\lift]_{\mathrm r})\\
 &=\PiTy(A[\rho;_{\Ren}\tau]_{\mathrm r},
          B[(\rho;_{\Ren}\tau)^\lift]_{\mathrm r}),
\end{aligned}
\]
using the two induction hypotheses and lift composition.
The same calculation covers \(\SigmaTy,\CodePi,\CodeSigma\); \(\Lam\)
uses only the bound induction hypothesis.  Application, pairing,
projections, successor, decoding, identity, code identity, and reflexivity
use the induction hypotheses componentwise.

For the two-binder constructors the recursive equations display all lifts:
\[
\begin{aligned}
&\NatElim_\ell(M,z,s,k)[\rho]_{\mathrm r}[\tau]_{\mathrm r}\\
&=\NatElim_\ell(
 M[\rho^\lift]_{\mathrm r}[\tau^\lift]_{\mathrm r},
 z[\rho]_{\mathrm r}[\tau]_{\mathrm r},
 s[\rho^{\lift\lift}]_{\mathrm r}
   [\tau^{\lift\lift}]_{\mathrm r},
 k[\rho]_{\mathrm r}[\tau]_{\mathrm r})\\
&=\NatElim_\ell(
 M[(\rho;_{\Ren}\tau)^\lift]_{\mathrm r},
 z[\rho;_{\Ren}\tau]_{\mathrm r},
 s[(\rho;_{\Ren}\tau)^{\lift\lift}]_{\mathrm r},
 k[\rho;_{\Ren}\tau]_{\mathrm r}),
\end{aligned}
\]
and
\[
\begin{aligned}
&\EqJ(M,r,b,p)[\rho]_{\mathrm r}[\tau]_{\mathrm r}\\
&=\EqJ(
 M[\rho^{\lift\lift}]_{\mathrm r}
   [\tau^{\lift\lift}]_{\mathrm r},
 r[\rho]_{\mathrm r}[\tau]_{\mathrm r},
 b[\rho]_{\mathrm r}[\tau]_{\mathrm r},
 p[\rho]_{\mathrm r}[\tau]_{\mathrm r})\\
&=\EqJ(
 M[(\rho;_{\Ren}\tau)^{\lift\lift}]_{\mathrm r},
 r[\rho;_{\Ren}\tau]_{\mathrm r},
 b[\rho;_{\Ren}\tau]_{\mathrm r},
 p[\rho;_{\Ren}\tau]_{\mathrm r}).
\end{aligned}
\]
The identity proof is the same induction with
\((\idren^{\Ren})^\lift=\idren^{\Ren}\), used twice in these final two
cases.  These are all constructors of \(\Expr\).
\end{proof}

\subsection{Simultaneous substitutions}

Set
\[
       \Sub(n,m)\coloneqq\Fin(n)\to\Expr_m.
\]
The embedding of renamings as variable-only substitutions is
\[
 \mathsf{ofRen}_{n,m}:\Ren(n,m)\to\Sub(n,m),\qquad
 \mathsf{ofRen}(\rho)(i)=\Var(\rho(i)).
\tag{4.2}\label{eq:of-ren}
\]
The substitution identity is not the renaming identity:
\[
       \idren^{\Sub}_n\coloneqq
       \mathsf{ofRen}(\idren^{\Ren}_n)\in\Sub(n,n).
\]
For \(\sigma\in\Sub(n,m)\), its lift lies in
\(\Sub(n+1,m+1)\) and is
\[
\sigma^\lift(0)=\Var(0),\qquad
\sigma^\lift(i+1)=\sigma(i)[\wk_m]_{\mathrm r}.
\tag{4.3}\label{eq:sub-lift}
\]
Capture-avoiding substitution \(t[\sigma]\) is again structural recursion,
using the same one- and two-lift pattern as renaming.  Composition is
\[
 \sigma;_{\Sub}\tau\in\Sub(n,k),\qquad
 (\sigma;_{\Sub}\tau)(i)=\sigma(i)[\tau],
\tag{4.4}\label{eq:sub-composition}
\]
for \(\sigma\in\Sub(n,m)\) and \(\tau\in\Sub(m,k)\).

The distinction between pre- and post-renaming is useful in the proof.
Writing \(\rho^\ast\sigma=\sigma\circ\rho\) and
\(\rho_\ast\sigma=(i\mapsto\sigma(i)[\rho]_{\mathrm r})\), one obtains the
fusion laws
\[
\begin{aligned}
t[\sigma][\rho]_{\mathrm r}&=t[\rho_\ast\sigma],\\
t[\rho]_{\mathrm r}[\sigma]&=t[\rho^\ast\sigma].
\end{aligned}
\tag{4.5}\label{eq:fusion}
\]
Their binder cases depend on the naturality square in
Proposition~\ref{prop:renaming}.

\begin{theorem}[Substitution algebra]
\label{thm:substitution-algebra}
For \(t\in\Expr_n\), \(\sigma\in\Sub(n,m)\), and
\(\tau\in\Sub(m,k)\), and \(\upsilon\in\Sub(k,l)\),
\[
\begin{aligned}
t[\idren^{\Sub}_n]&=t,\\
t[\sigma][\tau]&=t[\sigma;_{\Sub}\tau],\\
(\idren^{\Sub}_n)^\lift&=\idren^{\Sub}_{n+1},\\
(\sigma;_{\Sub}\tau)^\lift&=\sigma^\lift;_{\Sub}\tau^\lift.
\end{aligned}
\tag{4.6}\label{eq:sub-algebra}
\]
The substitution environments themselves satisfy
\[
\begin{aligned}
\idren^{\Sub}_n;_{\Sub}\sigma&=\sigma,&
\sigma;_{\Sub}\idren^{\Sub}_m&=\sigma,\\
(\sigma;_{\Sub}\tau);_{\Sub}\upsilon
 &=\sigma;_{\Sub}(\tau;_{\Sub}\upsilon).
\end{aligned}
\tag{4.7}\label{eq:sub-environment-algebra}
\]
Moreover
\[
                 t[\mathsf{ofRen}(\rho)]=t[\rho]_{\mathrm r},
\tag{4.8}\label{eq:ofren-subst}
\]
and the fusion equations~\eqref{eq:fusion} hold with their displayed scope
types.
\end{theorem}

\begin{proof}
We establish the lift and fusion equations before expression composition.
First,
\[
 \mathsf{ofRen}(\rho)^\lift=\mathsf{ofRen}(\rho^\lift)
\tag{4.9}\label{eq:ofren-lift}
\]
pointwise: both sides return \(\Var(0)\) at \(0\), and
\(\Var(\rho(i)+1)\) at \(i+1\).  Structural induction on \(t\), using
\eqref{eq:ofren-lift} once or twice under binders, proves
\eqref{eq:ofren-subst}.

Next prove the two fusion equations~\eqref{eq:fusion} by structural induction
on \(t\).  Their variable cases are the definitions of
\(\rho_\ast\sigma\) and \(\rho^\ast\sigma\).  Under one binder the required
environment equalities are, pointwise,
\[
\begin{aligned}
(\rho_\ast\sigma)^\lift(0)&=\Var(0),\\
(\rho_\ast\sigma)^\lift(i+1)
 &=\sigma(i)[\rho]_{\mathrm r}[\wk]_{\mathrm r}
  =\sigma(i)[\wk]_{\mathrm r}[\rho^\lift]_{\mathrm r},\\
(\rho^\ast\sigma)^\lift(0)&=\Var(0),\\
(\rho^\ast\sigma)^\lift(i+1)
 &=\sigma(\rho(i))[\wk]_{\mathrm r}
  =\sigma^\lift(\rho^\lift(i+1)).
\end{aligned}
\]
The middle equality in the first pair is the weakening naturality equation
of Proposition~\ref{prop:renaming}.  The two-binder constructors apply these
equalities twice.  All nonbinding constructors follow componentwise.

The identity lift is also pointwise:
\[
(\idren^{\Sub}_n)^\lift(0)=\Var(0),\qquad
(\idren^{\Sub}_n)^\lift(i+1)=\Var(i+1),
\]
so it is \(\idren^{\Sub}_{n+1}\).
For lift composition, both sides return \(\Var(0)\) at \(0\).  At \(i+1\),
\[
\begin{aligned}
(\sigma^\lift;_{\Sub}\tau^\lift)(i+1)
 &=\bigl(\sigma(i)[\wk_m]_{\mathrm r}\bigr)[\tau^\lift]\\
 &=\bigl(\sigma(i)[\tau]\bigr)[\wk_k]_{\mathrm r}\\
 &=(\sigma;_{\Sub}\tau)^\lift(i+1),
\end{aligned}
\tag{4.10}\label{eq:lift-composition-proof}
\]
where the middle step is the two fusion equations, equivalently the
imagewise square
\[
 \mathsf{ofRen}(\wk_m);_{\Sub}\tau^\lift
 =\tau;_{\Sub}\mathsf{ofRen}(\wk_k).
\]

We can now prove expression substitution.  Identity follows from
\eqref{eq:ofren-subst} and renaming identity.  For composition, induct
structurally on \(t\).  Variables give the definition of
\(\sigma;_{\Sub}\tau\); constants are fixed; all nonbinding constructors
use their induction hypotheses componentwise.  The one-binder cases use
\eqref{eq:lift-composition-proof}.  For the two-binder arguments,
\[
 M[\sigma^{\lift\lift}][\tau^{\lift\lift}]
  =M[\sigma^{\lift\lift};_{\Sub}\tau^{\lift\lift}]
  =M[(\sigma;_{\Sub}\tau)^{\lift\lift}],
\tag{4.11}\label{eq:eqj-double-lift-composition}
\]
and likewise for the successor branch \(s\) of \(\NatElim\).
Thus the \(\NatElim\) motive uses one lift, its successor branch two, and its
other arguments none; the \(\EqJ\) motive uses two lifts and its other
arguments none.  This exhausts the binding constructors.

Only after expression composition is established do we derive
\eqref{eq:sub-environment-algebra}.  At \(i\), left identity is
\(\Var(i)[\sigma]=\sigma(i)\), right identity is
\(\sigma(i)[\idren^{\Sub}_m]=\sigma(i)\), and associativity is
\[
 ((\sigma;_{\Sub}\tau);_{\Sub}\upsilon)(i)
 =(\sigma(i)[\tau])[\upsilon]
 =\sigma(i)[\tau;_{\Sub}\upsilon]
 =(\sigma;_{\Sub}(\tau;_{\Sub}\upsilon))(i).
\]
Function extensionality yields the environment equalities.  There is
therefore no use of environment associativity in the proof of expression
composition.
\end{proof}

\subsection{Weakening and instantiation}

For \(a\in\Expr_n\), let
\(\langle a\rangle\in\Sub(n+1,n)\) send \(0\) to \(a\) and \(i+1\) to
\(\Var(i)\).  One-variable instantiation is
\[
                     B\{a\}\coloneqq B[\langle a\rangle].
\]
For \(M\in\Expr_{n+2}\) and \(b,p\in\Expr_n\), let
\(\langle b,p\rangle\in\Sub(n+2,n)\) send \(0\) to \(p\), \(1\) to \(b\),
and \(i+2\) to \(\Var(i)\).  Then
\[
                     M\{b,p\}\coloneqq M[\langle b,p\rangle].
\]
The order reflects de Bruijn extension: the proof is newest.

\begin{proposition}[Instantiation laws]
\label{prop:instantiation}
Weakening is cancelled by one-variable instantiation,
\[
                   t[\wk]\{a\}=t,
\]
and by two-variable instantiation,
\[
              t[\wk][\wk]\{b,p\}=t.
\]
Renaming and substitution commute with both forms:
\[
\begin{aligned}
B\{a\}[\rho]_{\mathrm r}
 &=B[\rho^\lift]_{\mathrm r}\{a[\rho]_{\mathrm r}\},\\
B\{a\}[\sigma]
 &=B[\sigma^\lift]\{a[\sigma]\},\\
M\{b,p\}[\rho]_{\mathrm r}
 &=M[\rho^{\lift\lift}]_{\mathrm r}
       \{b[\rho]_{\mathrm r},p[\rho]_{\mathrm r}\},\\
M\{b,p\}[\sigma]
 &=M[\sigma^{\lift\lift}]\{b[\sigma],p[\sigma]\}.
\end{aligned}
\tag{4.12}\label{eq:instantiation-laws}
\]
\end{proposition}

\begin{proof}
Expand instantiation as simultaneous substitution and use expression
composition from Theorem~\ref{thm:substitution-algebra}.  For the
one-variable substitution law, the two induced substitutions in
\(\Sub(n+1,m)\) agree at \(0\), where both return \(a[\sigma]\), and at
\(i+1\), where both return \(\sigma(i)\).  Function extensionality and
expression composition give the displayed equality.

For the two-variable substitution law, both sides act on \(M\) by substitutions in
\(\Sub(n+2,m)\).  At index \(0\) both return \(p[\sigma]\); at index \(1\)
both return \(b[\sigma]\); at \(i+2\), one side gives
\(\Var(i)[\sigma]=\sigma(i)\), while the other gives
\(\Var(i+2)[\sigma^{\lift\lift}]\) and then instantiates, again reducing to
\(\sigma(i)\).  The renaming proof has the identical three cases with
\(\rho\) in place of \(\sigma\).  For weakening cancellation,
\(\mathsf{ofRen}(\wk_n);_{\Sub}\langle a\rangle\) is
\(\idren^{\Sub}_n\) pointwise, and the two-fold version satisfies
\[
\mathsf{ofRen}(\wk_n);_{\Sub}
\mathsf{ofRen}(\wk_{n+1});_{\Sub}\langle b,p\rangle
 =\idren^{\Sub}_n.
\]
Expression identity completes both equations.
\end{proof}

\section{Typing, computation, and primitive-redex preservation}
\label{sec:judgments}

\subsection{Typing rules}

The judgment \(\Gamma\vdash t:A\) is an inductive family, not a predicate
obtained from target semantics.  Variables and universes satisfy
\[
 \frac{i\in\Fin(n)}{\Gamma\vdash\Var(i):\Gamma(i)}
 \qquad
 \frac{}{\Gamma\vdash\Sort_\ell:\Sort_{\ell+1}}.
\]
The dependent product rules are
\[
\frac{\Gamma\vdash A:\Sort_\ell\qquad
      \Ext{\Gamma}{A}\vdash B:\Sort_\ell}
     {\Gamma\vdash\PiTy(A,B):\Sort_\ell}.
\]
\[
\frac{\Gamma\vdash A:\Sort_\ell\quad
      \Ext{\Gamma}{A}\vdash B:\Sort_\ell\quad
      \Ext{\Gamma}{A}\vdash t:B}
     {\Gamma\vdash\Lam(t):\PiTy(A,B)},
\]
\[
\frac{\Gamma\vdash f:\PiTy(A,B)\qquad\Gamma\vdash a:A}
     {\Gamma\vdash\App(f,a):B\{a\}}.
\]
The sum rules form \(\SigmaTy(A,B)\), introduce
\(\Pair(a,b)\) from \(a:A\) and \(b:B\{a\}\), and type the projections as
\[
  \Gamma\vdash\Fst(p):A,\qquad
  \Gamma\vdash\Snd(p):B\{\Fst(p)\}.
\]

For each \(\ell\), the natural-number rules are
\[
\Gamma\vdash\NatTy_\ell:\Sort_\ell,\quad
\Gamma\vdash\Zero_\ell:\NatTy_\ell,\quad
\frac{\Gamma\vdash k:\NatTy_\ell}
     {\Gamma\vdash\Succ_\ell(k):\NatTy_\ell}.
\]
If \(M\) is a type in \(\Ext{\Gamma}{\NatTy_\ell}\), the zero case has type
\(M\{\Zero_\ell\}\), and the successor case is checked in
\(\Ext{\Ext{\Gamma}{\NatTy_\ell}}{M}\) against the explicit substitution of
\(\Succ_\ell(\Var(1))\) for the natural argument, then the recursor has type
\(M\{k\}\).  Appendix~\ref{app:grammar} records the complete rule.

The universe-code rules include
\[
\frac{}{\Gamma\vdash\CodeNat_\ell:\Sort_\ell},
\qquad
\frac{\Gamma\vdash A:\Sort_\ell\quad
      \Ext{\Gamma}{\El(A)}\vdash B:\Sort_\ell}
     {\Gamma\vdash\CodePi(A,B):\Sort_\ell},
\]
with an analogous sum rule, and
\[
\frac{\Gamma\vdash A:\Sort_\ell\quad
      \Gamma\vdash a:\El(A)\quad\Gamma\vdash b:\El(A)}
     {\Gamma\vdash\CodeId(A,a,b):\Sort_\ell}.
\]
Any \(c:\Sort_\ell\) decodes to \(\El(c):\Sort_\ell\).

Identity formation and reflexivity are standard for the bounded grammar:
\[
\frac{\Gamma\vdash A:\Sort_\ell\quad\Gamma\vdash a:A\quad
      \Gamma\vdash b:A}
     {\Gamma\vdash\IdTy(A,a,b):\Sort_\ell},
\qquad
\frac{\Gamma\vdash a:A}
     {\Gamma\vdash\Refl(a):\IdTy(A,a,a)}.
\]
For raw \(\EqJ\), if \(M\) is a type expression in the motive
assignment~\eqref{eq:motive-context}, then
\[
\frac{
 \begin{array}{c}
 \Gamma\vdash A:\Sort_\ell\quad\Gamma\vdash a:A\\
 \Ext{\Ext{\Gamma}{A}}{\mathsf{IdProof}_n(A,a)}
      \vdash M:\Sort_j\\
 \Gamma\vdash r:M\{a,\Refl(a)\}\quad
 \Gamma\vdash b:A\quad
 \Gamma\vdash p:\IdTy(A,a,b)
 \end{array}}
 {\Gamma\vdash\EqJ(M,r,b,p):M\{b,p\}}.
\tag{5.1}\label{eq:eqj-rule}
\]

\begin{warning}[No conversion rule]
The raw judgment has no general rule converting a term from \(A\) to a
definitionally equal \(B\).  Nor is there a typed judgment
\(\Gamma\vdash A\defeq B\).  The preservation theorem below therefore
returns an explicit target type expression \(B\) and an \emph{untyped}
derivation \(B\defeq A\), rather than silently applying conversion.
\end{warning}

\subsection{Primitive computation and definitional equality}

The directed computation judgment \(r\red s\) has exactly the following
schemes:
\begin{equation}
\begin{aligned}
\App(\Lam(t),a)&\red t\{a\},\\
\Fst(\Pair(a,b))&\red a,
&
\Snd(\Pair(a,b))&\red b,\\
\NatElim_\ell(M,z,s,\Zero_\ell)&\red z,\\
\NatElim_\ell(M,z,s,\Succ_\ell(k))
 &\red s\{k,\NatElim_\ell(M,z,s,k)\},\\
\EqJ(M,r,a,\Refl(a))&\red r,\\
\El(\CodeNat_\ell)&\red\NatTy_\ell,\\
\El(\CodePi(A,B))&\red\PiTy(\El(A),\El(B)),\\
\El(\CodeSigma(A,B))&\red\SigmaTy(\El(A),\El(B)),\\
\El(\CodeId(A,a,b))&\red\IdTy(\El(A),a,b).
\end{aligned}
\tag{5.2}\label{eq:computations}
\end{equation}
Only raw expressions occur in these rules.  These are exactly ten
constructors of a top-level relation.  In particular, \(r\red s\) does not
mean reduction under an expression context, and it has no reflexive,
transitive, or multi-step constructor.  Contextual closure appears only in
the separate relation \(\defeq\).

Source definitional equality \(\defeq\) is generated by primitive
computation, reflexivity, symmetry, transitivity, and a constructor-by-
constructor compatible closure.  The closure is explicit: it covers both
arguments of products and sums, lambda bodies, function and argument
positions of application, pair components, projections, successor, all four
recursor positions, code arguments, decoding, the three identity positions,
reflexivity, and all four \(\EqJ\) positions.  This enumeration is important:
there is no clause saying that an arbitrary target trace creates source
definitional equality.

\begin{theorem}[Fundamental typing substitution]
\label{thm:fundamental-substitution}
Let \(\Gamma\in\Ctx(n)\), \(\Delta\in\Ctx(m)\), and
\(\sigma\in\Sub(n,m)\).  Suppose that \(\Gamma\vdash t:A\), and that
\(\sigma\) is typed from \(\Gamma\) to \(\Delta\), meaning
\[
      \Delta\vdash\sigma(i):\Gamma(i)[\sigma]
                 \quad\text{for every }i\in\Fin(n).
\tag{5.3}\label{eq:typed-substitution}
\]
Then
\[
                   \Delta\vdash t[\sigma]:A[\sigma].
\]
\end{theorem}

\begin{proof}
Induct on the typing derivation.  The variable case is
\eqref{eq:typed-substitution}.  The binder lemma says that a typed
\(\sigma\in\Sub(n,m)\) induces a typed
\(\sigma^\lift\in\Sub(n+1,m+1)\) from \(\Ext{\Gamma}{A}\) to
\(\Ext{\Delta}{A[\sigma]}\).  At index \(0\), the image is \(\Var(0)\), whose
declaration is \(A[\sigma][\wk_m]_{\mathrm r}\); this is also
\((A[\wk_n]_{\mathrm r})[\sigma^\lift]\) by the weakening square.  At
\(i+1\), the image \(\sigma(i)[\wk_m]_{\mathrm r}\) is typed by the renaming
admissibility theorem applied to
\(\Delta\vdash\sigma(i):\Gamma(i)[\sigma]\).  The type is rewritten by the
same square to
\((\Gamma(i)[\wk_n]_{\mathrm r})[\sigma^\lift]\).

We now inspect every typing constructor.  Universes, natural formation,
zero, and the natural code are unchanged.  Product and sum formation, lambda
introduction, and product/sum code formation apply the induction hypotheses
under \(\sigma^\lift\) using the binder lemma.  Application uses the
induction hypotheses for \(f\) and \(a\), then the one-variable equation
\[
 (B\{a\})[\sigma]=B[\sigma^\lift]\{a[\sigma]\}.
\]
Pairing uses the same equation for the second component.  First and second
projection, successor, decoding, identity/code-identity formation, and
reflexivity use their induction hypotheses componentwise; the second
projection again rewrites its displayed type by the one-variable equation.
This accounts for every constructor except \(\NatElim\) and \(\EqJ\).

For the natural successor branch, apply the binder lemma twice to
obtain a typed
\(\sigma^{\lift\lift}\in\Sub(n+2,m+2)\).  The induction hypothesis types the
substituted branch at
\(T_{\mathrm{succ}}(M)[\sigma^{\lift\lift}]\); the pointwise calculation
\[
 T_{\mathrm{succ}}(M)[\sigma^{\lift\lift}]
   =T_{\mathrm{succ}}(M[\sigma^\lift])
\tag{5.4}\label{eq:nat-target-substitution}
\]
has two cases on an index of \(M\): its natural argument \(0\), which is sent
to \(\Succ_\ell(\Var(1))\), and an inherited variable \(i+1\), shifted by
two.  The motive premise uses \(\sigma^\lift\), the zero case and scrutinee
use \(\sigma\), and the successor branch uses
\(\sigma^{\lift\lift}\).  The conclusion type is put in the required form by
\[
 M\{k\}[\sigma]=M[\sigma^\lift]\{k[\sigma]\}.
\]

The raw \(\EqJ\) motive uses the same double lift.  The identity-proof
declaration satisfies
\[
 \mathsf{IdProof}_n(A,a)[\sigma^\lift]
   =\mathsf{IdProof}_m(A[\sigma],a[\sigma]),
\tag{5.5}\label{eq:id-proof-substitution}
\]
because the weakened \(A\) and \(a\) commute with lift and \(\Var(0)\) is
fixed.  Thus the twice-lifted environment is typed for the substituted motive
assignment.  In the \(\EqJ\) case the induction hypotheses give, separately,
the substituted formation of \(A\), typing of \(a\), typing of
\(M[\sigma^{\lift\lift}]\) in that assignment, and typing of
\(r[\sigma]\), \(b[\sigma]\), and \(p[\sigma]\).  The proof premise has type
\[
 \IdTy(A[\sigma],a[\sigma],b[\sigma])
\]
by structural substitution.  The reflexivity-case type and conclusion type
are rewritten by the same two-variable law:
\[
 \begin{aligned}
 M\{a,\Refl(a)\}[\sigma]
  &=M[\sigma^{\lift\lift}]
       \{a[\sigma],\Refl(a[\sigma])\},\\
 M\{b,p\}[\sigma]
  =M[\sigma^{\lift\lift}]\{b[\sigma],p[\sigma]\},
 \end{aligned}
\]
the three-index calculation of Proposition~\ref{prop:instantiation}.  These
are exactly the premise and conclusion types of the substituted raw
\(\EqJ\) rule.
No quotient or target trace is used.
\end{proof}

\begin{corollary}[Typed instantiation]
\label{cor:typed-instantiation}
If \(\Ext{\Gamma}{A}\vdash t:T\) and \(\Gamma\vdash a:A\), then
\[
             \Gamma\vdash t\{a\}:T\{a\}.
\]
If \(\Ext{\Ext{\Gamma}{A}}{B}\vdash t:T\), \(\Gamma\vdash a:A\), and
\(\Gamma\vdash p:B\{a\}\), then
\[
             \Gamma\vdash t\{a,p\}:T\{a,p\}.
\]
\end{corollary}

\begin{theorem}[Computation and definitional equality under substitution]
\label{thm:computation-substitution}
For \(r,s\in\Expr_n\) and every \(\sigma\in\Sub(n,m)\),
\[
 r\red s\Longrightarrow r[\sigma]\red s[\sigma],
 \qquad
 r\defeq s\Longrightarrow r[\sigma]\defeq s[\sigma].
\]
\end{theorem}

\begin{proof}
Inspect the ten computation constructors.  Product beta is preserved because
\[
\App(\Lam(t),a)[\sigma]
 =\App(\Lam(t[\sigma^\lift]),a[\sigma])
 \red t[\sigma^\lift]\{a[\sigma]\}
 =(t\{a\})[\sigma].
\]
The two projection redexes become
\[
\Fst(\Pair(a[\sigma],b[\sigma]))\red a[\sigma],
\qquad
\Snd(\Pair(a[\sigma],b[\sigma]))\red b[\sigma].
\]
Natural-zero beta becomes
\[
\NatElim_\ell(M[\sigma^\lift],z[\sigma],
 s[\sigma^{\lift\lift}],\Zero_\ell)\red z[\sigma].
\]
Natural-successor beta is
\[
\begin{aligned}
&\NatElim_\ell(M[\sigma^\lift],z[\sigma],
 s[\sigma^{\lift\lift}],\Succ_\ell(k[\sigma]))\\
&\quad\red
s[\sigma^{\lift\lift}]
 \{k[\sigma],
   \NatElim_\ell(M[\sigma^\lift],z[\sigma],
       s[\sigma^{\lift\lift}],k[\sigma])\}\\
&\quad=
\bigl(s\{k,\NatElim_\ell(M,z,s,k)\}\bigr)[\sigma],
\end{aligned}
\]
by the two-variable instantiation equation.  Raw \(J\)-beta becomes
\[
\begin{aligned}
&\EqJ(M,r,a,\Refl(a))[\sigma]\\
&\quad=
\EqJ(M[\sigma^{\lift\lift}],r[\sigma],a[\sigma],
     \Refl(a[\sigma]))
\red r[\sigma],
\end{aligned}
\]
and the four decoding redexes become, respectively,
\[
\begin{gathered}
\El(\CodeNat_\ell)\red\NatTy_\ell,\\
\El(\CodePi(A[\sigma],B[\sigma^\lift]))
 \red\PiTy(\El(A[\sigma]),\El(B[\sigma^\lift])),\\
\El(\CodeSigma(A[\sigma],B[\sigma^\lift]))
 \red\SigmaTy(\El(A[\sigma]),\El(B[\sigma^\lift])),\\
\El(\CodeId(A[\sigma],a[\sigma],b[\sigma]))
 \red\IdTy(\El(A[\sigma]),a[\sigma],b[\sigma]).
\end{gathered}
\]
This exhausts the ten constructors.

For \(\defeq\), induct on its derivation.  Reflexivity, symmetry, and
transitivity use the corresponding closure constructors; a beta leaf uses
the preceding ten-case argument.  Product/sum domain congruence applies the
induction hypothesis under \(\sigma\), while codomain congruence applies it
under \(\sigma^\lift\); lambda and product/sum-code binders likewise use one
lift.  In \(\NatElim\), the motive uses \(\sigma^\lift\), the successor
branch uses \(\sigma^{\lift\lift}\), and the other two positions use
\(\sigma\).  In \(\EqJ\), the motive uses two lifts and the other three
positions use none.  Application, pairing, projections, successor, identity,
code identity, decoding, and reflexivity use the induction hypotheses in
their same-scope positions.  Reapplying the matching congruence constructor
proves each case, with no quotient assumption.
\end{proof}

Theorem~\ref{thm:computation-substitution} has one further consequence that is
used twice below: once in the redex-preservation theorem immediately following,
and again in Section~\ref{sec:syntactic-quotient} when substitution is
transported to the syntactic quotient.  We record it here, at the level of the
source relation, before any quotient has been formed.

\begin{proposition}[Substitution environment congruence]
\label{prop:environment-congruence}
Suppose \(\sigma,\tau\in\Sub(n,m)\) and
\(\sigma(i)\defeq\tau(i)\) for every \(i\in\Fin(n)\).  Then
\[
                     t[\sigma]\defeq t[\tau]
                     \quad\text{for every }t\in\Expr_n.
\tag{5.6}\label{eq:subst-congruence}
\]
\end{proposition}

\begin{proof}
Induct on \(t\).  The variable case is the assumption.  First prove the lift
lemma pointwise:
\[
 \sigma\equiv_{\mathrm{env}}\tau
 \Longrightarrow
 \sigma^\lift\equiv_{\mathrm{env}}\tau^\lift.
\tag{5.7}\label{eq:lift-env-congruence}
\]
At index \(0\), both sides are \(\Var(0)\).  At \(i+1\), apply
Theorem~\ref{thm:computation-substitution} to
\(\sigma(i)\defeq\tau(i)\) with
\(\mathsf{ofRen}(\wk_m)\); by~\eqref{eq:ofren-subst} this is precisely
\(\sigma(i)[\wk_m]_{\mathrm r}\defeq
  \tau(i)[\wk_m]_{\mathrm r}\).  Applying the lemma twice gives pointwise
\(\sigma^{\lift\lift}\equiv_{\mathrm{env}}
 \tau^{\lift\lift}\).

For \(\PiTy,\SigmaTy,\Lam,\CodePi,\CodeSigma\), use
\eqref{eq:lift-env-congruence} in the bound argument and the original
assumption in the unbound argument.  Application, pairing, projections,
successor, decoding, identity, code identity, and reflexivity use the
induction hypotheses in their same-scope positions, chaining one-position
congruence constructors when there is more than one argument.  Constants
are reflexive.  For \(\NatElim_\ell(M,z,s,k)\), the four induction hypotheses use,
respectively, \(\sigma^\lift\), \(\sigma\),
\(\sigma^{\lift\lift}\), and \(\sigma\); transitivity chains the four
one-position congruence constructors.  Concretely, it compares motives,
zero branches, successor branches, and scrutinees in that order:
\[
\begin{aligned}
&\NatElim_\ell(M[\sigma^\lift],z[\sigma],
 s[\sigma^{\lift\lift}],k[\sigma])\\
\defeq{}&
\NatElim_\ell(M[\tau^\lift],z[\sigma],
 s[\sigma^{\lift\lift}],k[\sigma])\\
\defeq{}&
\NatElim_\ell(M[\tau^\lift],z[\tau],
 s[\sigma^{\lift\lift}],k[\sigma])\\
\defeq{}&
\NatElim_\ell(M[\tau^\lift],z[\tau],
 s[\tau^{\lift\lift}],k[\sigma])\\
\defeq{}&
\NatElim_\ell(M[\tau^\lift],z[\tau],
 s[\tau^{\lift\lift}],k[\tau]).
\end{aligned}
\]
Each step changes exactly one displayed argument using the corresponding
congruence constructor.  Raw \(\EqJ(M,r,b,p)\) similarly uses
\(\sigma^{\lift\lift}\) for \(M\) and \(\sigma\) for the remaining three
arguments:
\[
\begin{aligned}
&\EqJ(M[\sigma^{\lift\lift}],r[\sigma],b[\sigma],p[\sigma])\\
\defeq{}&
\EqJ(M[\tau^{\lift\lift}],r[\sigma],b[\sigma],p[\sigma])\\
\defeq{}&
\EqJ(M[\tau^{\lift\lift}],r[\tau],b[\sigma],p[\sigma])\\
\defeq{}&
\EqJ(M[\tau^{\lift\lift}],r[\tau],b[\tau],p[\sigma])\\
\defeq{}&
\EqJ(M[\tau^{\lift\lift}],r[\tau],b[\tau],p[\tau]).
\end{aligned}
\]
Thus every constructor, including the two-binder cases, is covered by the
explicit compatible closure of \(\defeq\).
\end{proof}

\begin{theorem}[Primitive-redex preservation modulo type \(\defeq\)]
\label{thm:redex-preservation}
For \(\Gamma\in\Ctx(n)\) and \(r,s,A\in\Expr_n\), if
\(\Gamma\vdash r:A\) and \(r\red s\), then there is \(B\in\Expr_n\)
such that
\[
                  \Gamma\vdash s:B
                  \qquad\text{and}\qquad B\defeq A.
\tag{5.8}\label{eq:redex-preservation}
\]
Here \(B\defeq A\) is the untyped source relation applied to expressions
being used as types.  The theorem covers exactly the ten top-level schemes
in~\eqref{eq:computations}.  Contextual and multi-step reduction are treated
separately in Section~\ref{sec:reduction}, where the corresponding subject
reduction theorem is proved under the explicitly stated conversion hypothesis
of Definition~\ref{def:conversion-rules}.
\end{theorem}

\begin{proof}
Invert the typing derivation of the redex and inspect which of the ten
computation constructors was supplied.
The complete case analysis is summarized here; ``same'' means that the
target derivation has exactly the displayed source type.
\[
\begin{array}{c|c|c}
\text{redex}&\text{target term and chosen target type}&
  \text{type comparison}\\ \hline
\Pi\text{-}\beta&t\{a\}:B\{a\}&\text{same}\\
\Sigma\text{-fst}&a:A&\text{same}\\
\Sigma\text{-snd}&b:B\{a\}&
 B\{a\}\defeq B\{\Fst(\Pair(a,b))\}\\
\NatTy\text{-zero}&z:M\{\Zero_\ell\}&\text{same}\\
\NatTy\text{-succ}&s\{k,R_k\}:M\{\Succ_\ell(k)\}&\text{same}\\
\EqJ\text{-}\beta&r:M\{a,\Refl(a)\}&\text{same}\\
\El\text{-Nat}&\NatTy_\ell:\Sort_\ell&\text{same}\\
\El\text{-}\Pi&\PiTy(\El(A),\El(B)):\Sort_\ell&\text{same}\\
\El\text{-}\Sigma&\SigmaTy(\El(A),\El(B)):\Sort_\ell&\text{same}\\
\El\text{-Id}&\IdTy(\El(A),a,b):\Sort_\ell&\text{same}
\end{array}
\tag{5.9}\label{eq:redex-case-table}
\]

For product beta, inversion produces
\(\Ext{\Gamma}{A}\vdash t:B\) and \(\Gamma\vdash a:A\);
Corollary~\ref{cor:typed-instantiation} gives
\(\Gamma\vdash t\{a\}:B\{a\}\).
For first projection, inversion reaches the pair premise
\(\Gamma\vdash a:A\).  For natural-zero beta, inversion reaches the stored
zero-case premise \(\Gamma\vdash z:M\{\Zero_\ell\}\).

For the dependent second projection, inversion has the exact shape
\[
 \frac{\Gamma\vdash a:A\qquad\Gamma\vdash b:B\{a\}}
      {\Gamma\vdash\Snd(\Pair(a,b)):
         B\{\Fst(\Pair(a,b))\}}.
\]
The target \(b\) is already typed at \(B\{a\}\).  Primitive first-projection
beta gives
\(\Fst(\Pair(a,b))\defeq a\); symmetry gives
\(a\defeq\Fst(\Pair(a,b))\).  Compare the two substitutions in
\(\Sub(n+1,n)\):
\[
\begin{array}{lll}
\langle a\rangle(0)=a,
&\quad&
\langle\Fst(\Pair(a,b))\rangle(0)=\Fst(\Pair(a,b)),\\
\langle a\rangle(i+1)=\Var(i),
&&
\langle\Fst(\Pair(a,b))\rangle(i+1)=\Var(i).
\end{array}
\]
They are pointwise \(\defeq\)-equal by the symmetric beta step at \(0\) and
reflexivity at \(i+1\).  Applying substitution-environment congruence
(Proposition~\ref{prop:environment-congruence}) to the body \(B\) gives
\[
       B\{a\}\defeq B\{\Fst(\Pair(a,b))\}.
\]
This is why the theorem returns a possibly different type expression instead
of pretending that a conversion rule exists.

For natural successor beta, write
\[
 R_k=\NatElim_\ell(M,z,s,k).
\]
Inversion gives \(\Gamma\vdash k:\NatTy_\ell\), hence the recursor rule gives
\(\Gamma\vdash R_k:M\{k\}\).  The successor branch is typed in
\(\Ext{\Ext{\Gamma}{\NatTy_\ell}}{M}\) at
\(T_{\mathrm{succ}}(M)\).  Two-variable typed
instantiation with \(k\) and \(R_k\) yields
\[
 \Gamma\vdash s\{k,R_k\}:
     T_{\mathrm{succ}}(M)\{k,R_k\}.
\]
The pointwise substitution calculation
\[
 T_{\mathrm{succ}}(M)\{k,R_k\}=M\{\Succ_\ell(k)\}
\tag{5.10}\label{eq:nat-target-instantiation}
\]
is exact: the branch proof variable does not occur in
\(T_{\mathrm{succ}}(M)\), the predecessor becomes \(k\), and the motive
variable becomes \(\Succ_\ell(k)\).  This is the displayed source type of the
successor redex.

For raw \(J\)-beta, inversion reaches the stored premise
\(\Gamma\vdash r:M\{a,\Refl(a)\}\), exactly the redex type.

For \(\El(\CodeNat_\ell)\), inversion first yields
\(\Gamma\vdash\CodeNat_\ell:\Sort_\ell\); natural formation gives
\(\Gamma\vdash\NatTy_\ell:\Sort_\ell\).
For a product code, inversion yields
\[
\Gamma\vdash A:\Sort_\ell,\qquad
\Ext{\Gamma}{\El(A)}\vdash B:\Sort_\ell.
\]
Decoding formation gives the domain
\(\Gamma\vdash\El(A):\Sort_\ell\) and the codomain
\(\Ext{\Gamma}{\El(A)}\vdash\El(B):\Sort_\ell\); product formation therefore
types \(\PiTy(\El(A),\El(B))\) at \(\Sort_\ell\).
The sum-code case is identical with sum formation.  For an identity code,
inversion gives
\[
\Gamma\vdash A:\Sort_\ell,\quad
\Gamma\vdash a:\El(A),\quad
\Gamma\vdash b:\El(A);
\]
decoding formation and identity formation type
\(\IdTy(\El(A),a,b)\) at \(\Sort_\ell\).
All four source decoding redexes were displayed at \(\Sort_\ell\), so their
type-coherence proof is reflexivity.  This completes all ten cases without a
target-typing premise or a conversion rule.
\end{proof}

\section{Contextual reduction and conditional subject reduction}
\label{sec:reduction}

Theorem~\ref{thm:redex-preservation} concerns \emph{top-level} redexes.
It says nothing about reducing a subterm in place, nor about reducing more
than once.  This section supplies both closures and then determines exactly
what stops the corresponding subject-reduction theorem.

\subsection{The congruence and multi-step closures}

\begin{definition}[Contextual reduction]
\label{def:contextual-reduction}
Contextual one-step reduction \(t\ctxred u\) is generated by
\begin{itemize}[leftmargin=1.4em]
\item a head rule, embedding each primitive computation constructor; and
\item one congruence rule per raw constructor argument, mirroring the
  congruence closure of \(\defeq\) exactly, but without reflexivity, symmetry,
  or transitivity.
\end{itemize}
Multi-step reduction \(t\ctxredstar{}u\) is the reflexive-transitive closure of
\(\red\).
\end{definition}

Enumerating congruence by constructor rather than by an abstract
one-hole-context notion is deliberate.  It keeps the supported closure
auditable, in the same style as \(\defeq\) in
the definition of \(\defeq\) in Section~\ref{sec:judgments}, and it makes the following two proofs immediate
inductions with one case per rule.

\begin{proposition}[Soundness and substitution stability]
\label{prop:reduction-basic}
For all \(t,u\in\Expr_n\):
\begin{enumerate}[label=(\roman*)]
\item \(t\ctxred u\) implies \(t\defeq u\), and \(t\ctxredstar{}u\) implies
  \(t\defeq u\);
\item \(t\ctxred u\) implies \(\sigma(t)\red\sigma(u)\) for every
  \(\sigma:\Sub(n,m)\), and likewise for \(\ctxredstar\);
\item \(t\ctxredstar{}u\) implies \(\den t=\den u\) in \(\Q_n\).
\end{enumerate}
\end{proposition}

\begin{proof}
For (i), each head rule is a \(\defeq\)-atom by the beta constructor, and each
congruence rule of \(\ctxred\) has a congruence constructor of \(\defeq\) with
the same shape; the multi-step case chains these with transitivity.  For (ii),
the head case is stability of primitive computation under substitution, and
each congruence case applies the induction hypothesis under the appropriate
one- or two-variable lift.  Part (iii) is (i) followed by quotient soundness.
\end{proof}

\subsection{What blocks contextual subject reduction}

Consider reducing the argument of an application.  From
\(\Gamma\vdash\App(f,a):B\{a\}\) and \(a\ctxred a'\), inversion gives
\(\Gamma\vdash f:\PiTy(A,B)\) and \(\Gamma\vdash a:A\), and the induction
hypothesis gives \(\Gamma\vdash a':A\).  Reapplying the elimination rule yields
\[
 \Gamma\vdash\App(f,a'):B\{a'\},
\]
whereas the required conclusion is \(\Gamma\vdash\App(f,a'):B\{a\}\).  The two
displayed types are related by \(B\{a\}\defeq B\{a'\}\), which follows from
\(a\defeq a'\) by congruence of instantiation.  Passing between them is
exactly a conversion step, and the raw typing family has no such rule.

A second obstruction appears under binders.  Reducing a domain annotation,
\(A\ctxred A'\) inside \(\PiTy(A,B)\), leaves the codomain premise typed in the
context \(\Ext{\Gamma}{A}\) while the reassembled term needs it in
\(\Ext{\Gamma}{A'}\).  The two contexts are pointwise \(\defeq\)-equal.
Passing between them is a context-conversion step, which the raw system also
omits.

These are the only two obstructions arising in the induction.  Before naming
them we settle the status of the theorem they block, because a blocked proof and
a false statement are different things.

\subsection{Contextual subject reduction is false}
\label{subsec:sr-fails}

The obstruction above is real, not an artifact of a clumsy induction.

\begin{theorem}[Failure of contextual subject reduction]
\label{thm:sr-fails}
Let \(\Gamma_0\) be the one-slot context assigning
\(\PiTy(\NatTy_0,\Var_0)\), let
\[
 r=\App(\Lam(\Var_0),\Zero_0)\ \ctxred\ \Zero_0,
 \qquad
 t=\App(\Var_0,r),
 \qquad
 u=\App(\Var_0,\Zero_0).
\tag{6.1}\label{eq:sr-counterexample}
\]
Then \(t\ctxred u\) and \(\Gamma_0\vdash t:r\) are derivable, while
\(\Gamma_0\vdash u:r\) is not.
\end{theorem}

\begin{proof}
The codomain of \(\Gamma_0(0)\) is the bound variable, so
\(\Var_0\{a\}=a\) and the application rule gives \(\Gamma_0\vdash t:r\) from the
variable rule and \(\Gamma_0\vdash r:\NatTy_0\); the latter is one lambda
introduction followed by one application.  The step \(t\ctxred u\) is the
argument congruence applied to the \(\Pi\)-beta redex.

For the negative half, suppose \(\Gamma_0\vdash u:T\).  The only rule with
\(\App\) in its conclusion is the application rule, so
\(T=B\{\Zero_0\}\) for some \(B\) with \(\Gamma_0\vdash\Var_0:\PiTy(A,B)\).  The
only rule with a variable in its conclusion reads the type off the context slot
verbatim, so \(\PiTy(A,B)=\Gamma_0(0)\) and hence \(B=\Var_0\).  Therefore
\(T=\Zero_0\).  Since \(r\) and \(\Zero_0\) have different head constructors,
\(T\neq r\).
\end{proof}

Two features of the raw calculus make the argument this short.  Contexts are
arbitrary same-scope assignments, so no well-formedness obligation is incurred
by \(\Gamma_0\); and the typing rules are syntax-directed, so the negative half
is inversion rather than a normalization argument.

Rather than leave the two obstructions as a limitation, we name them.

\begin{definition}[The omitted structural rules]
\label{def:conversion-rules}
The \emph{conversion hypothesis} consists of two rules:
\[
\frac{\Gamma\vdash t:A\qquad A\defeq A'}{\Gamma\vdash t:A'}
\qquad
\frac{\Gamma\vdash t:A\qquad \Gamma\defeq\Delta}
     {\Delta\vdash t:A}
\tag{6.2}\label{eq:conversion-rules}
\]
where \(\Gamma\defeq\Delta\) abbreviates pointwise source definitional
equality of the two same-scope assignments.
\end{definition}

The hypothesis is assumed as data, never as an axiom.  Every statement below
carries it as an explicit argument, so its dependence on the omitted rules is
visible in the statement itself.

It is visible for a reason: the hypothesis cannot be met.

\begin{theorem}[Neither rule is admissible]
\label{thm:not-admissible}
For the raw judgment \(\vdash\):
\begin{enumerate}[label=(\roman*)]
\item conversion is not admissible.  In any context,
  \(\vdash\Zero_0:\NatTy_0\) is derivable and
  \(\NatTy_0\defeq\El(\CodeNat_0)\), yet
  \(\vdash\Zero_0:\El(\CodeNat_0)\) is not derivable;
\item context conversion is not admissible.  The family
  \(\IdTy(\El(\CodeNat_0),\Var_0,\Var_0)\) is a type over the domain
  \(\El(\CodeNat_0)\) and is not a type over the pointwise
  \(\defeq\)-equal domain \(\NatTy_0\);
\item consequently the record of Definition~\ref{def:conversion-rules} is
  uninhabited.
\end{enumerate}
\end{theorem}

\begin{proof}
For (i), the only rule concluding a type for \(\Zero_0\) assigns
\(\NatTy_0\) exactly, and \(\NatTy_0\) and \(\El(\CodeNat_0)\) are distinct raw
expressions; the definitional equality is one instance of the decoding
computation rule.  For (ii), inverting the identity-formation rule leaves a
premise typing \(\Var_0\) at \(\El(\CodeNat_0)\), and the variable rule reads
that type off the context slot verbatim, forcing the weakened slot to be
\(\El(\CodeNat_0)\); over the reduced domain the slot is \(\NatTy_0\) instead.
That the family \emph{is} a type over the unreduced domain is one application of
identity formation, so the failure is a genuine loss rather than an artifact of
an ill-formed example.  Part (iii) follows from (i) alone.
\end{proof}

Theorem~\ref{thm:not-admissible} is the necessity counterpart of
Theorem~\ref{thm:sr-fails}, and it forces us to be careful about how the
conditional theorem below is read.  Two congruence lemmas are needed to state
it.

\begin{lemma}[Context and instantiation congruence]
\label{lem:ctx-congruence}
Pointwise \(\defeq\) of contexts is preserved by extension, hence also by the
natural-successor context \(\mathsf{natStepCtx}\) and the equality-motive
context of Section~\ref{sec:judgments}.  Moreover, instantiation is congruent
in its \emph{body}: \(M\defeq M'\) implies \(M\{a\}\defeq M'\{a\}\) and
\(M\{a,p\}\defeq M'\{a,p\}\).
\end{lemma}

\begin{proof}
For extension, both the new declaration and every weakened old declaration are
obtained by applying the same renaming to \(\defeq\)-equal expressions, and
\(\defeq\) is stable under substitution.  The two derived contexts are
extensions.  Body congruence is stability of \(\defeq\) under substitution
applied to the instantiating environment, which is fixed.
\end{proof}

\begin{theorem}[Contextual subject reduction]
\label{thm:contextual-subject-reduction}
Assume the conversion hypothesis \eqref{eq:conversion-rules}.  If
\(\Gamma\vdash t:A\) and \(t\ctxred u\), then \(\Gamma\vdash u:A\), at the same
type \(A\).
\end{theorem}

\begin{proof}
Induction on the derivation of \(t\ctxred u\).

The head case is Theorem~\ref{thm:redex-preservation}, which supplies a
target typing at a type \(\defeq\)-equal to \(A\); the conversion rule moves it
to \(A\).

Each congruence case inverts the syntax-directed typing rule for the
surrounding constructor, applies the induction hypothesis to the reduced
premise, and reassembles.  When the surrounding rule does not mention the
reduced subterm in its conclusion type---codomains, lambda bodies,
\(\Sigma\)-second components, the \(\mathsf{natElim}\) branches, decoding, and
the identity endpoints---reassembly is immediate.

When it does, conversion repairs the displayed type.  For
\(\App(f,a)\) with \(a\ctxred a'\), the reassembled type \(B\{a'\}\) is moved to
\(B\{a\}\) using \(a\defeq a'\) and congruence of instantiation.  For
\(\Snd(p)\) with \(p\ctxred p'\), the same argument uses
\(\Fst(p)\defeq\Fst(p')\).  For \(\mathsf{natElim}\) with a reduced
scrutinee, and for \(\EqJ\) with a reduced endpoint or proof, the argument is
identical with the corresponding instantiation congruence.  For a reduced
\(\mathsf{natElim}\) motive or \(\EqJ\) motive, body congruence of
Lemma~\ref{lem:ctx-congruence} is used instead, once on the branch premises and
once on the conclusion.  For \(\Refl(a)\) with \(a\ctxred a'\), both endpoints of
the identity type change and are repaired together.

When the reduced subterm is a domain annotation, context conversion repairs
the premise: for \(\PiTy(A,B)\) with \(A\ctxred A'\) the codomain premise is
moved from \(\Ext{\Gamma}{A}\) to \(\Ext{\Gamma}{A'}\) by
Lemma~\ref{lem:ctx-congruence}, and similarly for \(\SigmaTy\), for the two
universe codes \(\CodePi\) and \(\CodeSigma\) with \(\El(A)\defeq\El(A')\),
and for the \(\mathsf{natElim}\) motive.
\end{proof}

Because the paper insists elsewhere that its constructor enumerations are
auditable, Table~\ref{tab:sr-cases} lists every congruence case of the
induction and the repair it needs, including the dependent-premise cases that
the prose above compresses.

\begin{table}[t]
\centering
\small
\renewcommand{\arraystretch}{1.18}
\begin{tabular}{>{\raggedright\arraybackslash}p{0.30\textwidth}
                >{\raggedright\arraybackslash}p{0.28\textwidth}
                >{\raggedright\arraybackslash}p{0.32\textwidth}}
\toprule
Reduced position & Conclusion type & Repair needed\\
\midrule
\(\Pi\), \(\Sigma\) codomain; \(\Lam\) body; \(\CodePi\), \(\CodeSigma\)
 codomain
 & \(\Sort_\ell\) or unchanged
 & none\\
\(\Pi\), \(\Sigma\) domain
 & \(\Sort_\ell\)
 & context conversion on the codomain premise\\
\(\CodePi\), \(\CodeSigma\) domain
 & \(\Sort_\ell\)
 & context conversion, via \(\El(A)\defeq\El(A')\)\\
\(\App\) function
 & \(B\{a\}\), unchanged
 & none\\
\(\App\) argument
 & \(B\{a\}\to B\{a'\}\)
 & conversion, instantiation congruence\\
\(\Pair\) first component
 & \(\SigmaTy(A,B)\), unchanged
 & none, but the second premise is retyped by conversion at
   \(B\{a\}\defeq B\{a'\}\)\\
\(\Pair\) second component; \(\Fst\)
 & unchanged
 & none\\
\(\Snd\)
 & \(B\{\Fst(p)\}\to B\{\Fst(p')\}\)
 & conversion, via \(\Fst(p)\defeq\Fst(p')\)\\
\(\Succ\) argument; \(\El\) code; identity endpoints
 & unchanged
 & none\\
\(\IdTy\) domain
 & \(\Sort_\ell\)
 & conversion on both endpoint premises\\
\(\Refl\) argument
 & \(\IdTy(A,a,a)\to\IdTy(A,a',a')\)
 & conversion, both endpoints together\\
\(\mathsf{natElim}\) zero or successor branch
 & \(M\{t\}\), unchanged
 & none\\
\(\mathsf{natElim}\) motive
 & \(M\{t\}\to M'\{t\}\)
 & conversion on the conclusion and on both branch premises; context
   conversion for the successor branch\\
\(\mathsf{natElim}\) scrutinee
 & \(M\{t\}\to M\{t'\}\)
 & conversion, instantiation congruence\\
\(\EqJ\) reflexivity case
 & unchanged
 & none\\
\(\EqJ\) motive
 & \(M\{b,p\}\to M'\{b,p\}\)
 & conversion on the conclusion and on the reflexivity-case premise\\
\(\EqJ\) endpoint or proof
 & \(M\{b,p\}\to M\{b',p\}\) etc.
 & conversion, two-variable instantiation congruence\\
\bottomrule
\end{tabular}
\caption{Every congruence case of
Theorem~\ref{thm:contextual-subject-reduction} and the repair it requires.  Only
the two shapes of \eqref{eq:conversion-rules} occur.}
\label{tab:sr-cases}
\end{table}

\begin{corollary}[Multi-step subject reduction]
\label{cor:multistep-subject-reduction}
Assume \eqref{eq:conversion-rules}.  If \(\Gamma\vdash t:A\) and
\(t\ctxredstar{}u\), then \(\Gamma\vdash u:A\).
\end{corollary}

\begin{proof}
Induction on the length of the reduction, applying
Theorem~\ref{thm:contextual-subject-reduction} at each step.
\end{proof}

\begin{remark}[How to read the conditional form]
\label{rem:conditional-value}
Theorem~\ref{thm:contextual-subject-reduction} is not a claim that the raw
calculus enjoys subject reduction; by Theorem~\ref{thm:sr-fails} it does not.
Nor is its hypothesis one a reader could hope to discharge for \(\vdash\): by
Theorem~\ref{thm:not-admissible} that hypothesis is unsatisfiable, so as a
statement \emph{about} \(\vdash\) the theorem is vacuously true.  We state it
anyway, and state this caveat with it, because its proof carries information
that neither of the negative theorems does: it enumerates every case of the
induction and shows that repairs of exactly the two shapes in
\eqref{eq:conversion-rules}---and no others---suffice to close all of them.
This is the sense, and the only sense, in which the pair is claimed to be
sufficient; it is a fact about the proof, not a measurement of the distance
between two calculi.
\end{remark}

\subsection{The conversion-closed extension}
\label{subsec:conversion-closed}

The constructive response to Theorem~\ref{thm:not-admissible} is to change the
judgment rather than to assume rules about it.

\begin{definition}[Conversion-closed judgment]
\label{def:vdashC}
Let \(\vdashC\) be generated by the raw typing rules together with the two
rules of \eqref{eq:conversion-rules}, read now as constructors rather than as
hypotheses.
\end{definition}

\begin{proposition}[Basic properties of \(\vdashC\)]
\label{prop:vdashC}
\(\vdashC\) contains \(\vdash\); it admits conversion and context conversion by
construction; and it types the reduct that defeated \(\vdash\) in
Theorem~\ref{thm:sr-fails}, namely \(\Gamma_0\vdashC u:r\).
\end{proposition}

\begin{proof}
The first two are immediate from the constructors.  For the third, the
application rule gives \(\Gamma_0\vdashC u:\Zero_0\) and conversion moves the
type along \(\Zero_0\defeq r\), which is the \(\Pi\)-beta rule read backwards.
\end{proof}

We stop there deliberately.

\begin{remark}[Why no unconditional theorem for \(\vdashC\)]
\label{rem:no-unconditional}
Proving contextual subject reduction for \(\vdashC\) outright would require
inverting \(\vdashC\), and \(\vdashC\) is not syntax-directed: a derivation may
end in conversion.  Inversion therefore has to be replaced by generation
lemmas, which for an \emph{untyped} definitional equality such as \(\defeq\)
require a confluence analysis; Remark~\ref{rem:reduction-limits} records that
we have none.  So the honest summary of this section is not that the raw
calculus is two rules away from subject reduction, but that it is a
\emph{different judgment} away, and that the metatheory of that judgment needs
confluence.  Supplying it is the natural next step, and it is exactly the point
at which this development stops being elementary.
\end{remark}

\begin{remark}[What is still not proved]
\label{rem:reduction-limits}
No confluence, normalization, or standardization theorem is claimed for
\(\red\).  Proposition~\ref{prop:reduction-basic} gives soundness and
substitution stability only.  In particular the syntactic quotient is
constructed from \(\defeq\) directly, not from a normal-form analysis of
\(\red\).
\end{remark}

\section{The syntactic quotient and compositional substitution}
\label{sec:syntactic-quotient}

\subsection{The quotient carrier}

For each scope \(n\), source definitional equality is a setoid on
\(\Expr_n\).  Define
\[
              \Q_n=\Expr_n/{\defeq},
       \qquad \eta_n(t)=\den t.
\tag{7.1}\label{eq:syntactic-quotient}
\]
This is only a syntactic quotient at a fixed scope.  It carries neither a
context-formation judgment nor the structure of a CwF, and no initiality
property is asserted.  No normal forms are chosen.

For \(\sigma\in\Sub(n,m)\), put
\[
       \Q(\sigma)(\den t)=\den{t[\sigma]}.
\tag{7.2}\label{eq:quotient-substitution}
\]
Theorem~\ref{thm:computation-substitution} proves that this definition is
independent of the representative \(t\).  This order matters logically:
substitution closure is a theorem about the source relation, proved before
the quotient action is defined.

\begin{corollary}[Representative independence in the environment]
\label{cor:quotient-environment}
If \(\sigma,\tau\in\Sub(n,m)\) satisfy \(\sigma(i)\defeq\tau(i)\) for every
\(i\in\Fin(n)\), then \(\Q(\sigma)=\Q(\tau)\) pointwise, and as maps under
function extensionality.
\end{corollary}

\begin{proof}
Proposition~\ref{prop:environment-congruence} gives \(t[\sigma]\defeq t[\tau]\)
for every \(t\in\Expr_n\), so quotient soundness gives
\(\Q(\sigma)(\den t)=\Q(\tau)(\den t)\), and quotient induction extends this to
an arbitrary \(x:\Q_n\).
\end{proof}

The action~\eqref{eq:quotient-substitution} is therefore independent of the
expression representative, by Theorem~\ref{thm:computation-substitution}, and of
the environment representative, by Corollary~\ref{cor:quotient-environment}.

\begin{theorem}[Compositionality of quotient substitution]
\label{thm:quotient-compositionality}
For every \(x\in\Q_n\) and composable substitutions
\(\sigma\in\Sub(n,m)\), \(\tau\in\Sub(m,k)\),
\[
\begin{aligned}
\Q(\idren^{\Sub}_n)(x)&=x,\\
\Q(\tau)(\Q(\sigma)(x))&=\Q(\sigma;_{\Sub}\tau)(x).
\end{aligned}
\tag{7.3}\label{eq:quotient-compositionality}
\]
The equations hold for arbitrary quotient values, not merely for a chosen
notation of the form \(\den t\).  If \(\sigma\) and \(\tau\) are pointwise
definitionally equal, their induced maps agree.
\end{theorem}

\begin{proof}
Define \(\Q(\sigma)\) by quotient elimination from
\[
 t\longmapsto\den{t[\sigma]}.
\]
Its representative-independence obligation is: from \(t\defeq u\), prove
\(\den{t[\sigma]}=\den{u[\sigma]}\).  Theorem~\ref{thm:computation-substitution}
gives \(t[\sigma]\defeq u[\sigma]\), and quotient soundness supplies exactly
that equality.  Thus the map is well defined.

For identity and composition, use quotient induction to reduce an arbitrary
\(x\) to \(\den t\).  Equations~\eqref{eq:sub-algebra} give
\(\den{t[\idren^{\Sub}_n]}=\den t\) and
\[
 \begin{aligned}
 \Q(\idren^{\Sub}_n)(\den t)
   &=\den{t[\idren^{\Sub}_n]}=\den t,\\
 \Q(\tau)(\Q(\sigma)(\den t))
   &=\den{t[\sigma][\tau]}
     =\den{t[\sigma;_{\Sub}\tau]}
     =\Q(\sigma;_{\Sub}\tau)(\den t).
 \end{aligned}
\]
Corollary~\ref{cor:quotient-environment} proves that pointwise
\(\defeq\)-equal environments induce pointwise equal quotient maps; function
extensionality gives equality of maps.  Hence the construction is independent
both of expression representatives and of pointwise equivalent environment
representatives, and the two displayed equations hold on every quotient
value.
\end{proof}

Theorem~\ref{thm:quotient-compositionality} asserts only these identity and
composition equations.  We have not defined a category of well-formed
contexts and substitutions, so no categorical structure is asserted.

\subsection{Action on quoted expressions}

For \(t\in\Expr_n\) and \(\sigma\in\Sub(n,m)\), abbreviate
\[
                       \den t_\sigma\coloneqq
                       \Q(\sigma)(\den t)=\den{t[\sigma]}.
\tag{7.4}\label{eq:environment-interpretation}
\]
Then
\[
          \den{t[\sigma]}_\tau
            =\den t_{\sigma;_{\Sub}\tau}.
\tag{7.5}\label{eq:interpretation-composition}
\]
The quotient equality in~\eqref{eq:interpretation-composition} may be
decorated by a singleton unlabelled quotient trace.  That decoration does
not remember the substitution proof and plays no role in establishing the
equation.

Typing can be saturated by quotient representatives without calling the
quotient a typed model.

\begin{definition}[Quotient-saturated typing]
\label{def:quotient-typing}
For \(x,X\in\Q_n\), write
\(\Gamma\models_{\Q}x:X\) when there exist \(t,A\in\Expr_n\) such that
\[
             \Gamma\vdash t:A,\qquad x=\den t,\qquad X=\den A.
\]
\end{definition}

\begin{proposition}[Substitution for quotient-saturated typing]
\label{prop:quotient-typing-substitution}
If \(\Gamma\models_{\Q}x:X\) and \(\sigma\) is a typed substitution from
\(\Gamma\) to \(\Delta\), then
\[
        \Delta\models_{\Q}\Q(\sigma)(x):\Q(\sigma)(X).
\]
\end{proposition}

\begin{proof}
Choose the representatives in Definition~\ref{def:quotient-typing} and
apply Theorem~\ref{thm:fundamental-substitution}.  The equalities identifying
\(x\) and \(X\) with their representatives are preserved by the function
\(\Q(\sigma)\).
\end{proof}

This relation merely packages a raw representative and two quotient
equalities.  It supplies no context validity, conversion, or initial CwF.

\section{Source identity programs and unlabelled quotient traces}
\label{sec:identity-programs}

\subsection{Finite same-scope frames and source programs}

For \(t\in\Expr_n\), let \(C[t]\) denote plugging \(t\) into the one-hole
frame \(C\).  The finite frame grammar used here is
\[
\begin{array}{llll}
\App([-],a),&\App(f,[-]),&
\Pair([-],b),&\Pair(a,[-]),\\
\Fst([-]),&\Snd([-]),&\Succ_\ell([-]),&\El([-]),\\
\Refl([-]),&
\IdTy([-],a,b),&\IdTy(A,[-],b),&\IdTy(A,a,[-]),\\
\CodeId([-],a,b),&\CodeId(A,[-],b),&\CodeId(A,a,[-]).
\end{array}
\tag{8.1}\label{eq:frames}
\]
These are same-scope frames; they do not include binder contexts.

\begin{lemma}[Frame congruence]
\label{lem:frame-congruence}
If \(t\defeq u\), then \(C[t]\defeq C[u]\) for every
\(C\) in~\eqref{eq:frames}.  Hence
\[
 \overline C:\Q_n\to\Q_n,\qquad
 \overline C(\den t)=\den{C[t]}
\]
is representative-independent.
\end{lemma}

\begin{proof}
Each case is the corresponding one-position congruence constructor of
\(\defeq\).  Quotient soundness proves representative independence.
\end{proof}

\begin{definition}[Source identity programs]
\label{def:identity-programs}
For \(t,u\in\Expr_n\), the family \(\IProg_n(t,u)\) is generated by
\[
\frac{h:t\defeq u}{\mathsf{atom}(h):\IProg_n(t,u)},\qquad
\mathbf{1}_t:\IProg_n(t,t),
\]
\[
\frac{p:\IProg_n(t,u)}{p^{-1}:\IProg_n(u,t)},\qquad
\frac{p:\IProg_n(t,u)\quad q:\IProg_n(u,v)}
     {p\cdot q:\IProg_n(t,v)}.
\]
\[
\frac{p:\IProg_n(t,u)}
     {C[p]:\IProg_n(C[t],C[u])}.
\]
\end{definition}

The atom accepts an untyped source \(\defeq\)-derivation.  It cannot accept a
target trace step or target trace equivalence.  The family itself has no
typing judgment for its endpoints.

\begin{proposition}[Endpoint triviality]
\label{prop:iprogram-endpoints}
For \(t,u\in\Expr_n\),
\[
 \operatorname{Nonempty}\bigl(\IProg_n(t,u)\bigr)
 \quad\Longleftrightarrow\quad t\defeq u.
\tag{8.2}\label{eq:iprogram-endpoint-iff}
\]
\end{proposition}

\begin{proof}
For the forward implication, choose \(p:\IProg_n(t,u)\) and induct on
\(p\).  An atom contains the required derivation.  Unit uses reflexivity;
inverse uses symmetry; composition uses transitivity; and \(C[p]\) uses
Lemma~\ref{lem:frame-congruence}.  Conversely, from \(h:t\defeq u\) take
\(\mathsf{atom}(h):\IProg_n(t,u)\).
\end{proof}

Thus identity-program syntax adds compositional trace structure only after
endpoint reachability has already been supplied by the untyped relation
\(\defeq\).  It cannot discover a new pair of reachable endpoints, and its
terms need not relate expressions having a common raw type.

\begin{proposition}[Structural-recursion interpretation]
\label{prop:iprogram-recursion}
Fix \(n\).  Let \(G(t,u)\) be any family indexed by
\(t,u\in\Expr_n\), equipped with:
\[
\begin{array}{ll}
a_h:G(t,u)&\text{for every }h:t\defeq u,\\
e_t:G(t,t),&\\
i_{t,u}:G(t,u)\to G(u,t),&\\
m_{t,u,v}:G(t,u)\to G(u,v)\to G(t,v),&\\
F_C:G(t,u)\to G(C[t],C[u])&\text{for every }C\in\Frame_n.
\end{array}
\tag{8.3}\label{eq:target-signature}
\]
There is a unique pointwise family of maps
\[
        \Phi_{t,u}:\IProg_n(t,u)\to G(t,u)
\tag{8.4}\label{eq:iprogram-fold}
\]
such that
\[
\begin{aligned}
\Phi(\mathsf{atom}(h))&=a_h,&
\Phi(\mathbf1_t)&=e_t,\\
\Phi(p^{-1})&=i(\Phi p),&
\Phi(p\cdot q)&=m(\Phi p)(\Phi q),\\
\Phi(C[p])&=F_C(\Phi p).
\end{aligned}
\tag{8.5}\label{eq:iprogram-fold-equations}
\]
\end{proposition}

\begin{proof}
Existence is structural recursion on the five constructors of
\(\IProg\).  If \(\Psi\) satisfies the same equations, induction on
\(p:\IProg_n(t,u)\) gives \(\Phi(p)=\Psi(p)\) in each constructor case.
This is pointwise uniqueness; function extensionality converts it to equality
of the map families if desired.
\end{proof}

Proposition~\ref{prop:iprogram-recursion} is a free-\emph{syntax} property
only.  It imposes no groupoid equations on \(G\) and says nothing about the
quotient of \(\IProg\) by source rewrites.  A universal property for a free
groupoid quotient would require proving that all source relations are
respected and is not claimed here.

\subsection{The unlabelled quotient-trace record}

\begin{definition}[Trace steps and trace records]
\label{def:trace-record}
For a metatheoretic carrier \(X\), define
\[
 \mathsf{QStep}(X)\coloneqq\sum_{x,y:X}(x=y),
 \qquad
 \Tr_X(a,b)\coloneqq
   \operatorname{List}(\mathsf{QStep}(X))\times(a=b).
\tag{8.6}\label{eq:trace-record}
\]
If \(S=(L,e):\Tr_X(a,b)\), write
\(\operatorname{list}(S)=L\) and
\(\operatorname{end}(S)=e\).
\end{definition}

The basic operations are fully determined on both fields.  Put
\[
\begin{aligned}
\mathbf r_a&=([],\mathsf{refl}_a),\\
(L,e)^{-1}
  &=\bigl(\operatorname{map}(\mathsf{inv},
       \operatorname{reverse}(L)),e^{-1}\bigr),\\
(L,e)\cdot(M,d)&=(L{+\!\!+}M,e\mathbin{\cdot_{=}}d),
\end{aligned}
\tag{8.7}\label{eq:trace-operations}
\]
where
\(\mathsf{inv}(x,y,h)=(y,x,h^{-1})\), \(+\!\!+\) is list
concatenation, and \(\cdot_{=}\) is transitivity of ambient equality.
For \(f:X\to Y\),
\[
 f_\ast(L,e)=
 \bigl(\operatorname{map}(x,y,h\mapsto
       (fx,fy,\mathsf{ap}_f(h)),L),\mathsf{ap}_f(e)\bigr).
\tag{8.8}\label{eq:trace-map}
\]
Thus \(f_\ast\mathbf r_a=\mathbf r_{fa}\),
\(f_\ast(S^{-1})=(f_\ast S)^{-1}\), and
\(f_\ast(S\cdot T)=f_\ast S\cdot f_\ast T\).

For a completely arbitrary record, Definition~\ref{def:trace-record} imposes
no adjacency relation between successive list entries and no relation between
the list and the outer endpoint field.  This fact is not used as a positive
invariant.  The soundness theorem quantifies only over expressions of the
form \(\ev(p)\).

When \(X=\Q_n\), the record is \emph{unlabelled}.  Given
\(h:t\defeq u\), quotient soundness produces
\(\bar h:\den t=\den u\), and the atom is evaluated by
\[
 \langle\bar h\rangle
 =\bigl([(\den t,\den u,\bar h)],\bar h\bigr).
\]
The derivation \(h\) is absent from this record.  In particular, two
\(\defeq\)-derivations with the same endpoints yield equal quotient-equality
proofs by proof irrelevance.  This differs from the label-sensitive calculus
of~\cite{deQueirozEtAl2016,RamosEtAl2017,RamosEtAl2018,RamosEtAl2026}.
Because \(\den t=\den u\) already in \(\Q_n\), an evaluated identity program
is best regarded as a trace-decorated loop over that quotient-collapsed
endpoint (although the displayed record retains the two representatives).

Evaluation is structural:
\[
\begin{aligned}
\ev(\mathsf{atom}(h))&=\langle\bar h\rangle,&
\ev(\mathbf{1}_t)&=\mathbf r_{\den t},\\
\ev(p^{-1})&=\ev(p)^{-1},&
\ev(p\cdot q)&=\ev(p)\cdot\ev(q),\\
\ev(C[p])&=\overline C_\ast(\ev(p)).
\end{aligned}
\tag{8.9}\label{eq:evaluation}
\]
This is the instance of Proposition~\ref{prop:iprogram-recursion} with
\[
 G(t,u)=\Tr_{\Q_n}(\den t,\den u),\quad
 a_h=\langle\bar h\rangle,\quad e_t=\mathbf r_{\den t},
\]
and with inverse, composition, and frame actions supplied by
\eqref{eq:trace-operations} and~\eqref{eq:trace-map}.
If desired, the evaluated sublanguage can be named without asserting an
unproved trace invariant:
\[
 \mathsf{EvalTr}_n(t,u)
   =\{\,T:\Tr_{\Q_n}(\den t,\den u)
        \mid \exists p:\IProg_n(t,u),\ T=\ev(p)\,\}.
\tag{8.10}\label{eq:evaluated-image}
\]
We use only the witness \(p\) and the defining equations~\eqref{eq:evaluation}.

\subsection{The target trace rewrite relation}

We now give the entire target relation needed below.  At fixed endpoints,
the one-step relation \(S\TrStep T\) has exactly eight schemes:
\[
\begin{array}{rclcrcl}
(\mathbf r_a)^{-1}&\TrStep&\mathbf r_a,
&& (S^{-1})^{-1}&\TrStep&S,\\
\mathbf r_a\cdot S&\TrStep&S,
&&S\cdot\mathbf r_b&\TrStep&S,\\
S\cdot S^{-1}&\TrStep&\mathbf r_a,
&&S^{-1}\cdot S&\TrStep&\mathbf r_b,\\
(S\cdot T)^{-1}&\TrStep&T^{-1}\cdot S^{-1},
&&(S\cdot T)\cdot U&\TrStep&S\cdot(T\cdot U).
\end{array}
\tag{8.11}\label{eq:target-one-step}
\]
Here \(S:\Tr_X(a,b)\), \(T:\Tr_X(b,c)\), and
\(U:\Tr_X(c,d)\) where required.  These rules are equations of trace
expressions oriented from left to right; they do not inspect source syntax.

\begin{definition}[Target trace equivalence]
\label{def:target-trace-equivalence}
\(\TrEq\) is the least endpoint-indexed family generated by
\[
\frac{}{S\TrEq S},\qquad
\frac{S\TrStep T}{S\TrEq T},\qquad
\frac{S\TrEq T}{T\TrEq S},\qquad
\frac{S\TrEq T\quad T\TrEq U}{S\TrEq U},
\tag{8.12}\label{eq:target-equivalence-closure}
\]
and the congruence rules
\[
\frac{S\TrEq T}{S^{-1}\TrEq T^{-1}},\qquad
\frac{S\TrEq T}{S\cdot U\TrEq T\cdot U},\qquad
\frac{T\TrEq U}{S\cdot T\TrEq S\cdot U},\qquad
\frac{S\TrEq T}{f_\ast S\TrEq f_\ast T}.
\tag{8.13}\label{eq:target-congruence}
\]
\end{definition}

Equations~\eqref{eq:target-one-step}--\eqref{eq:target-congruence} are a
complete mathematical presentation of the target fragment used in the
soundness proof.  They are the groupoid core and congruence fragment of the
larger computational-path rewrite system, not a presentation of all rules in
the monograph.

\subsection{Source rewrites and soundness}

The source one-step relation \(p\srw q\) has the eight schemes obtained from
\eqref{eq:target-one-step} by replacing \(S,T,U,\mathbf r\) with
\(p,q,r,\mathbf1\).  It also has the congruence rules
\[
\frac{p\srw q}{p^{-1}\srw q^{-1}},\quad
\frac{p\srw q}{p\cdot r\srw q\cdot r},\quad
\frac{q\srw r}{p\cdot q\srw p\cdot r},\quad
\frac{p\srw q}{C[p]\srw C[q]},
\tag{8.14}\label{eq:source-congruence-rules}
\]
and the frame-unit rule
\[
                  C[\mathbf1_t]\srw\mathbf1_{C[t]}.
\tag{8.15}\label{eq:frame-unit}
\]
Let \(\srweq\) be the reflexive, symmetric, transitive closure of
\(\srw\).

\begin{theorem}[Source identity rewrite soundness]
\label{thm:identity-soundness}
For \(p,q:\IProg_n(t,u)\),
\[
          p\srweq q
          \quad\Longrightarrow\quad
          \ev(p)\TrEq\ev(q).
\tag{8.16}\label{eq:rewrite-soundness}
\]
\end{theorem}

\begin{proof}
Induct first on a source one-step derivation.  The eight groupoid cases are
the eight target rules~\eqref{eq:target-one-step} after unfolding
\eqref{eq:evaluation}.  Inverse and the two composition congruences use the
first three rules in~\eqref{eq:target-congruence}.  Frame congruence uses the
last rule with \(f=\overline C\).  The frame-unit case is reflexivity because
\(\overline C_\ast\mathbf r_{\den t}=\mathbf r_{\den{C[t]}}\) by
\eqref{eq:trace-map}.  Induct next on the reflexive, symmetric, transitive
closure of the source relation and use the first, third, and fourth rules in
\eqref{eq:target-equivalence-closure}.  No target derivation is an input to a
source constructor.
\end{proof}

\begin{example}[Multi-step soundness]
\label{ex:two-step-coherence}
For \(p:\IProg_n(t,u)\) and \(q:\IProg_n(u,v)\), source associativity
followed by right-unit congruence gives
\[
       (p\cdot q)\cdot\mathbf{1}_v
          \srweq p\cdot q.
\]
The soundness derivation first applies target associativity and then removes
the right unit inside the right factor.  This is a two-step derivation of
trace equivalence; no equality between derivations is asserted.
\end{example}

\section{Derivation erasure: the traces are unlabelled}
\label{sec:erasure}

Section~\ref{sec:identity-programs} called the evaluated traces
\emph{unlabelled}, on the grounds that a source \(\defeq\)-derivation is used
only to obtain a proof-irrelevant equality of quotient classes.  Stated that
way it is a remark about intent.  Since the computational-path literature
adopts the opposite convention---the equality reason is a first-class
label---the difference is worth a theorem, and this section supplies one.

The strategy is to exhibit a syntax that manifestly \emph{cannot} carry a
label, and to show that evaluation already goes through it.

\subsection{A label-free program syntax}

\begin{definition}[Quotient programs]
\label{def:quot-programs}
For \(a,b:\Q_n\), the type \(\mathsf{QProg}_n(a,b)\) is generated by
\[
\begin{aligned}
 \mathsf{atom} &: (a=b)\to\mathsf{QProg}_n(a,b), \\
 \mathbf 1_a &: \mathsf{QProg}_n(a,a), \\
 (-)^{-1} &: \mathsf{QProg}_n(a,b)\to\mathsf{QProg}_n(b,a), \\
 (-)\cdot(-) &: \mathsf{QProg}_n(a,b)\to\mathsf{QProg}_n(b,c)
   \to\mathsf{QProg}_n(a,c), \\
 f_{*} &: \mathsf{QProg}_n(a,b)\to\mathsf{QProg}_n(f(a),f(b)),
\end{aligned}
\tag{9.1}\label{eq:quot-programs}
\]
the last for any \(f:\Q_n\to\Q_n\).  Its evaluation
\(\ev_{\Q}:\mathsf{QProg}_n(a,b)\to\Tr_{\Q_n}(a,b)\) sends \(\mathsf{atom}\,e\)
to the singleton trace of \(e\) and the remaining constructors to the
corresponding trace operations.
\end{definition}

The point of \eqref{eq:quot-programs} is what it lacks.  Its atoms accept an
equality in \(\Q_n\), which is a proof-irrelevant proposition.  There is
therefore no position in a quotient program at which a source derivation could
be stored, inspected, or branched on: the type has no room for one.  Any
statement proved about \(\mathsf{QProg}\) is automatically a statement that
does not depend on source derivations.

\begin{definition}[Erasure]
\label{def:erasure}
Erasure \(\lfloor-\rfloor:\IProg_n(t,u)\to\mathsf{QProg}_n(\den t,\den u)\) is
defined by structural recursion:
\[
\begin{aligned}
 \lfloor\mathsf{atom}\,h\rfloor &=\mathsf{atom}\,(\mathsf{sound}(h)),
 &\quad
 \lfloor\mathbf 1_t\rfloor &=\mathbf 1_{\den t}, \\
 \lfloor p^{-1}\rfloor &=\lfloor p\rfloor^{-1},
 &\quad
 \lfloor p\cdot q\rfloor &=\lfloor p\rfloor\cdot\lfloor q\rfloor, \\
 \lfloor C_{*}p\rfloor &=(\Q C)_{*}\lfloor p\rfloor,
\end{aligned}
\tag{9.2}\label{eq:erasure-clauses}
\]
where \(\mathsf{sound}\) is quotient soundness and \(\Q C\) is the quotient
action of the frame \(C\) from Section~\ref{sec:identity-programs}.
\end{definition}

The atom clause is the whole content: a derivation \(h\) of \(t\defeq u\) is
replaced by the single equality \(\den t=\den u\) that it justifies.  Because
that equality is proof-irrelevant, the replacement is irreversible in the sense
that \(h\) cannot be read back off.  Everything else is a relabelling of
constructors.

\subsection{Evaluation factors through erasure}

\begin{theorem}[Evaluation factors through erasure]
\label{thm:erasure-factorization}
For every \(p:\IProg_n(t,u)\),
\[
 \ev(p)=\ev_{\Q}\bigl(\lfloor p\rfloor\bigr)
 \qquad\text{in }\Tr_{\Q_n}(\den t,\den u).
\tag{9.3}\label{eq:erasure-factorization}
\]
\end{theorem}

\begin{proof}
Structural induction on \(p\).  For an atom, both sides are by definition the
singleton trace of the same quotient equality; the two proofs of that equality
are equal by proof irrelevance, so the traces are equal.  Reflexivity is the
empty trace on both sides.  The inverse, composition, and frame cases are the
induction hypotheses transported through the corresponding trace operation,
which is a congruence.
\end{proof}

\begin{corollary}[No target observation sees a derivation]
\label{cor:no-derivation-observation}
If \(h_1,h_2:t\defeq u\) are two source derivations of the same definitional
equality, then \(\ev(\mathsf{atom}\,h_1)=\ev(\mathsf{atom}\,h_2)\).  More
generally, no function defined on evaluated traces can separate two source
programs with equal erasures.
\end{corollary}

\begin{proof}
Immediate from Theorem~\ref{thm:erasure-factorization}: the two programs have
the same erasure, hence equal evaluations, hence equal images under any
function.  For the atom case this is already proof irrelevance of
\(\den t=\den u\).
\end{proof}

Corollary~\ref{cor:no-derivation-observation} is exactly the failure of
label-sensitivity.  In a label-carrying calculus the analogous statement is
false by design: distinct reasons for the same equality remain distinguishable
downstream, which is what makes their rewrite theory informative.

\subsection{Shape is preserved, and atoms round trip}

Two further clauses are needed before the picture is complete.  Erasure could
in principle be uninformative for the wrong reason---by discarding structure as
well as labels---or too generous, by admitting quotient atoms that no source
derivation supports.  Neither happens.

\begin{proposition}[Size preservation]
\label{prop:erasure-size}
Erasure preserves the number of program nodes:
\(|\lfloor p\rfloor|=|p|\) for every \(p:\IProg_n(t,u)\).
\end{proposition}

\begin{proof}
Induction on \(p\); each clause of \eqref{eq:erasure-clauses} maps a
constructor to a constructor of the same arity.
\end{proof}

So the arrangement of equality steps survives erasure intact.  What is
discarded is precisely the labels, and Proposition~\ref{prop:erasure-size}
separates this from the degenerate situation in which a trace would forget its
own shape.

\begin{proposition}[Atom correspondence and round trip]
\label{prop:erasure-roundtrip}
For \(t,u\in\Expr_n\):
\begin{enumerate}[label=(\roman*)]
\item \(t\defeq u\) holds if and only if \(\den t=\den u\);
\item writing \(\mathsf{exact}\) for the recovery map of quotient exactness,
  \[
   \bigl\lfloor\mathsf{atom}\,(\mathsf{exact}(e))\bigr\rfloor
   =\mathsf{atom}\,e
   \qquad\text{for every }e:\den t=\den u.
  \tag{9.4}\label{eq:atom-roundtrip}
  \]
\end{enumerate}
\end{proposition}

\begin{proof}
Part (i) is quotient soundness in one direction and quotient exactness in the
other.  For (ii), both sides are \(\mathsf{atom}\) applied to an equality in
\(\Q_n\), and any two such equalities between the same classes are equal by
proof irrelevance.
\end{proof}

Part (ii) says a quotient atom can always be realized by a source derivation,
and that erasing that derivation returns the atom one started from.  This is an
\emph{atom-level} statement, and we resist the temptation to read it as a
global one: it does not say that source programs and label-free programs are
equivalent languages.  Indeed they are not, and the next subsection says
exactly how far apart they are.

\subsection{The essential image of erasure}
\label{subsec:erasure-image}

The label-free grammar \eqref{eq:quot-programs} is deliberately permissive.  Its
congruence constructor accepts an arbitrary endofunction \(f:\Q_n\to\Q_n\),
whereas the source grammar offers only the finite collection of frames of
Section~\ref{sec:identity-programs}, and erasure supplies only \(\Q\)-actions of
the form \(C[-]\).  Erasure is therefore \emph{not} surjective onto
\(\QProg\), and a factorization theorem alone does not identify the two
languages.  What can be identified is the image.

\begin{definition}[Framed programs]
\label{def:framed}
Say that \(q:\QProg_n(\den t,\den u)\) is \emph{framed}, written
\(\Framed(t,u,q)\), when it is generated by:
atoms \(\mathsf{atom}(\mathsf{sound}(h))\) for \(h:t\defeq u\); reflexivity
\(\mathbf1_{\den t}\); inverses and composites of framed programs; and
congruences \(q\mapsto C[q]\) for a source frame \(C\).
\end{definition}

\begin{theorem}[Essential image]
\label{thm:erasure-image}
For all \(t,u\in\Expr_n\) and \(q:\QProg_n(\den t,\den u)\),
\[
 \bigl(\exists\,p:\IProg_n(t,u).\ \lfloor p\rfloor=q\bigr)
 \quad\Longleftrightarrow\quad
 \Framed(t,u,q).
\tag{9.5}\label{eq:erasure-image}
\]
\end{theorem}

\begin{proof}
Left to right is induction on \(p\): each erasure clause produces the
corresponding framed constructor, atoms going to quotient soundness of a source
derivation and congruences to \(C[-]\) for the frame that \(p\) supplied.  Right
to left is induction on the derivation of \(\Framed\), reading each clause
backwards to build the source program; the frame clause is where the restriction
is used, since an unframed congruence has no source preimage to offer.
\end{proof}

Two consequences are worth stating separately, because they are the precise
versions of two slogans we avoid.  Erasure \emph{invents no congruence}: every
congruence it produces is \(C[-]\) for an actual source frame, never an
arbitrary endofunction.  And erasure \emph{invents no atom}: every atom it
produces is quotient soundness applied to a source definitional equality.  Both
are immediate from the erasure clauses, and together with
Theorem~\ref{thm:erasure-image} they replace the unsupported claim that erasure
``identifies the target exactly''.

\begin{corollary}[Rewrite soundness through erasure]
\label{cor:erased-soundness}
Source rewrite soundness passes to erased programs: if \(p\srweq q\) then
\(\ev_{\Q}(\lfloor p\rfloor)\TrEq\ev_{\Q}(\lfloor q\rfloor)\).
\end{corollary}

\begin{proof}
Rewrite \eqref{eq:erasure-factorization} on both sides of
Theorem~\ref{thm:identity-soundness}.
\end{proof}

Corollary~\ref{cor:erased-soundness} completes the separation of layers begun
in Section~\ref{sec:layers}: the entire target theory---evaluation, traces, and
their rewrite equivalence---can be developed over the label-free syntax, and
the source calculus contributes only through which atoms exist.

\begin{remark}[Relation to label-sensitive computational paths]
\label{rem:erasure-vs-labels}
The results of this section make the comparison of
Section~\ref{sec:comparison} precise rather than impressionistic.  Earlier
computational-path work retains the equality reason and studies its rewrite
calculus; there, an analogue of
Corollary~\ref{cor:no-derivation-observation} would be false, and that failure
is the source of the theory's content.  Here the reason is erased by
construction, and Theorem~\ref{thm:erasure-factorization} shows the erasure is
already implicit in evaluation.  Neither target subsumes the other: the present
one supports quotient-level reasoning and proof irrelevance, the label-sensitive
one supports intensional distinctions this target cannot express.
\end{remark}

\section{Examples: beta, dependency, and two-binder instantiation}
\label{sec:example}

\subsection{Lambda beta across the quotient layers}

Fix a level \(\ell\) and work in the empty scope.  In one bound variable let
\(x=\Var(0)\), and consider
\[
       r=\App(\Lam(x),\Zero_\ell),
       \qquad s=\Zero_\ell.
\tag{10.1}\label{eq:running-redex}
\]
The variable rule gives
\(\Ext{\EmptyCtx_0}{\NatTy_\ell}\vdash
x:\NatTy_\ell[\wk_0]_{\mathrm r}\); the latter expression is
definitionally the constant \(\NatTy_\ell\).  Product formation and lambda
introduction therefore derive
\[
  \EmptyCtx_0\vdash
     \Lam(x):\PiTy(\NatTy_\ell,\NatTy_\ell),
\]
where the codomain notation suppresses its one-variable scope.  Together
with
\(\EmptyCtx_0\vdash\Zero_\ell:\NatTy_\ell\), application gives
\[
                  \EmptyCtx_0\vdash r:\NatTy_\ell.
\]

At the operational layer, product beta gives
\[
          r\red x\{\Zero_\ell\}=\Zero_\ell=s.
\tag{10.2}\label{eq:running-beta}
\]
The final equality is the variable clause of simultaneous substitution.
Embedding primitive computation into the untyped definitional relation gives
\[
                         r\defeq s.
\tag{10.3}\label{eq:running-defeq}
\]
Consequently quotient soundness identifies their classes:
\[
                         \den r=\den s
                          \quad\text{in }\Q_0.
\tag{10.4}\label{eq:running-quotient}
\]

Let \(\beta\) denote the derivation~\eqref{eq:running-defeq} and define the
source identity atom
\[
                   \alpha=\mathsf{atom}(\beta)
                     :\IProg_0(r,s).
\]
Its evaluation is the singleton unlabelled quotient trace
\[
 P_\beta=\ev(\alpha)
  =\Bigl(
      [\,(\den r,\den s,\overline\beta)\,],
      \overline\beta
    \Bigr)
 :\Tr_{\Q_0}(\den r,\den s).
\tag{10.5}\label{eq:running-path}
\]
It has one visible list entry, but that entry contains only the quotient
equality \(\overline\beta\), not the derivation \(\beta\).  Source
cancellation is sound as
\[
    \alpha\cdot\alpha^{-1}\srw\mathbf{1}_r
    \quad\longmapsto\quad
    P_\beta\cdot P_\beta^{-1}
        \TrEq\mathbf r_{\den r}.
\tag{10.6}\label{eq:running-rweq}
\]

Thus the example exhibits source beta, untyped source definitional equality,
quotient equality, an unlabelled trace decoration, and target trace
equivalence.  The quotient step erases the computational label; no arrow is
reversed to define a source judgment.

The same calculation is natural in environments.  If
\(\App(\Lam(t),a)\in\Expr_n\) and \(\sigma\in\Sub(n,m)\), then
\[
\App(\Lam(t),a)[\sigma]
   \red t[\sigma^\lift]\{a[\sigma]\}
   =t\{a\}[\sigma].
\]
The final equation is the one-variable case of
Proposition~\ref{prop:instantiation}.

\subsection{Dependent projection and natural successor, in brief}
\label{sec:example-snd}

The remaining examples follow the same three-layer pattern and we state only
what is specific to each.  For the dependent second projection, from
\(\Gamma\vdash a:A\) and \(\Gamma\vdash b:B\{a\}\) the pairing rule gives
\(\Gamma\vdash\Pair(a,b):\SigmaTy(A,B)\), and the second-projection redex
\(\Snd(\Pair(a,b))\red b\) is typed on both sides at \(B\{\Fst(\Pair(a,b))\}\)
and \(B\{a\}\) respectively.  These two type expressions are \(\defeq\)-equal
but not identical, which is exactly the situation
Theorem~\ref{thm:redex-preservation} is stated to accommodate and exactly the
situation that defeats subject reduction on the nose.

For the natural successor, the branch derivation lives under two binders,
\(\Ext{\Ext{\Gamma}{\NatTy_\ell}}{M}\vdash s:T_{\mathrm{succ}}(M)\).  Typed
two-variable instantiation at \(\Gamma\vdash k:\NatTy_\ell\) and
\(\Gamma\vdash R_k:M\{k\}\) yields
\(\Gamma\vdash s\{k,R_k\}:T_{\mathrm{succ}}(M)\{k,R_k\}\), and the explicit
substitution defining \(T_{\mathrm{succ}}\) gives
\(T_{\mathrm{succ}}(M)\{k,R_k\}=M\{\Succ_\ell(k)\}\), which is the type of the
redex.  The equality is a three-index calculation on the substitution into
\(M\), not an appeal to target traces.

\subsection{\texorpdfstring{What raw \(J\)-beta says}{What raw J-beta says}}

If the six premises of~\eqref{eq:eqj-rule} specialize to endpoint \(a\) and
proof \(\Refl(a)\), then
\[
             \EqJ(M,r,a,\Refl(a))\red r
\]
and the source and target are typed at the same raw expression
\(M\{a,\Refl(a)\}\).  This is the entire raw result.  It does not say that
\(M\) factors through quotient equality and is not interpreted by the
external construction of Section~\ref{sec:j-boundary}.

\section{\texorpdfstring{What the traces are: the metadata-fiber criterion}
{What the traces are: the metadata-fiber criterion}}
\label{sec:j-boundary}

It remains to say what the objects produced by Section~\ref{sec:erasure}
actually are.  The answer is deflationary, and stating it plainly is the point
of this section: a trace is \emph{observable list metadata decorating an
equality}, not a carrier of reasons for it.

Three facts, all already proved above, force that reading.  A trace is a list of
equality steps over \(\Q_n\) together with an endpoint equality; entries are not
required to be adjacent, composable, or related to the outer endpoints
(Definition~\ref{def:trace-record}).  Each step is an equality in a
proof-irrelevant set, so it marks a pair of already-collapsed classes rather
than a reduction that produced one from the other.  And the source derivation
that justified a step is not merely unused but provably unrecoverable
(Theorem~\ref{thm:erasure-factorization},
Corollary~\ref{cor:no-derivation-observation}).  A reader who expects a trace
calculus to expose reduction histories should read those results as saying that
the histories have been erased.

The companion article \cite{ObservableMetadata2026} classifies exactly when such
a decoration still supports identity elimination.  We use its criterion; we do
not redevelop it.

\begin{theorem}[Metadata-fiber criterion, {\cite{ObservableMetadata2026}}]
\label{thm:metadata-fiber-criterion}
Fix \(a:A\), a family \(M:\prod_{b:A}(a=b)\to\mathcal U\), and
\(m_0:M(a,\mathsf{refl}_a)\).  The refinement
\(R_M(a,b)=\sum_{h:a=b}M(b,h)\), based at \((\mathsf{refl}_a,m_0)\), admits an
unrestricted based eliminator into all \(\mathcal U\)-valued motives with
propositional beta if and only if \(M(a,\mathsf{refl}_a)\) is contractible.
\end{theorem}

Applying it to our own records is immediate once the record is put in the shape
of a refinement, which is a reordering of fields.

\begin{theorem}[Trace instance]
\label{thm:no-trace-j}
For every \(X\) and \(a:X\), the trace record \(\Tr_X\) with reflexivity
\(\mathbf r\) does not admit an unrestricted based eliminator satisfying
propositional beta.
\end{theorem}

\begin{proof}
Take \(M_X(b,h)=\operatorname{List}(\mathsf{QStep}(X))\).  The maps
\((L,e)\mapsto(e,L)\) and \((h,L)\mapsto(L,h)\) are mutually inverse by
reflexivity of ordered pairs, so \(\Tr_X(a,b)\simeq\sum_{h:a=b}M_X(b,h)\).  The
reflexivity fiber is \(\operatorname{List}(\mathsf{QStep}(X))\), which contains
the empty list and the singleton
\([\,(a,a,\mathsf{refl}_a)\,]\); these are distinct by no-confusion, so the fiber is not
contractible and Theorem~\ref{thm:metadata-fiber-criterion} applies.
\end{proof}

The singleton used here is well formed even under the usual intrinsic chain
conditions: its sole step starts and ends at \(a\), its outer endpoints are
\(a,a\), and its composite equality is reflexivity.  Imposing composability
therefore does not repair the obstruction, and neither does any intrinsic
metadata family retaining distinct empty and singleton reflexive values.  What
remains available is elimination for motives factoring through the
endpoint-equality projection
\[
 \Tr_X(a,b)\longrightarrow(a=b),\qquad (L,e)\longmapsto e,
\tag{11.1}\label{eq:forget-trace}
\]
which cannot observe the list; this is the projection theorem of
\cite{ObservableMetadata2026}, and it is what the present traces do support.

\section{Boundaries and possible strengthenings}
\label{sec:limitations}

This section is the consolidated list of what the audit does \emph{not}
establish.  It is stated once, here, and the remaining sections do not restate
it.

\paragraph{Source and reduction.}
The source has only the constructors in Section~\ref{sec:syntax}.  Contexts
are same-scope assignments, with no context-formation judgment, typed
\(\defeq\), or conversion rule.  The paper proves substitution and
preservation for ten top-level redexes unconditionally, proves that contextual
subject reduction fails, and proves the conversion hypothesis of
Definition~\ref{def:conversion-rules} unsatisfiable.  Contextual and
multi-step subject reduction are obtained only under that hypothesis and are
therefore vacuous as statements about \(\vdash\); no unconditional
subject-reduction theorem is proved for \(\vdashC\) either, because we prove
no normalization, confluence, standardization, canonicity, or decidable type
checking.

\paragraph{Identity programs and endpoints.}
\(\IProg\) is untyped and endpoint-trivial.  We prove only the homomorphic
simulation \(\srweq\Rightarrow\TrEq\), with no converse, completeness, or
claim that every trace is evaluated.  Distinct empty and singleton metadata
does not refute UIP for ambient equality; it shows that the metadata fiber is
noncontractible.  The singleton is intrinsically composable, so
composability alone cannot produce an unrestricted based eliminator
satisfying propositional beta.

\paragraph{Quotient and external algebra.}
\(\Tr_X(a,b)\) presupposes \(a=b\), and endpoint exactness is immediate from
the chosen quotient.  No interpretation of the raw calculus is supplied; one
would have to cover assignments, levels, codes, typing, and all redexes.
Univalence, higher inductive types, and a pentagon equality between rewrite
derivations are not claimed.

\paragraph{Erasure.}
Theorem~\ref{thm:erasure-factorization} says evaluation factors through the
label-free syntax; it does not say the two syntaxes are isomorphic.  The
quotient syntax admits frame actions along arbitrary endomaps of \(\Q_n\),
whereas source programs use only the finite frame grammar, so erasure is not
surjective.  Theorem~\ref{thm:erasure-image} identifies the essential image
instead, and Proposition~\ref{prop:erasure-roundtrip} establishes the exact
correspondence at atoms; nothing stronger is asserted.

\section{Comparison with prior computational-path results}
\label{sec:comparison}

Table~\ref{tab:comparison} compares statements, not implementation size or
historical importance.  ``Not stated in this form'' is deliberately weaker
than a priority claim: the table records how the cited theorems are used and
how the present statements differ.

The comparison draws on identifiable portions of the sources: normal forms,
labelled deduction, and groupoid structure in \cite{deQueirozEtAl2016};
computational paths, the path-based identity construction, and groupoid laws in
\cite{RamosEtAl2017}; the path term-rewriting system and its type-theoretic
constructions in \cite{RamosEtAl2018}; and labelled natural deduction, the path
rewrite system, the groupoid and \(\omega\)-groupoid structures, and \(J\),
UIP and transport in the 2026 monograph \cite{RamosEtAl2026}.  Those are all
label-sensitive developments.  This article instead uses an unlabelled
quotient-trace target as a source-separation experiment.  We do not claim that
those works lack substitution results; the claim is only that the same-scope
raw interface assembled here is isolated and proved as a single package.

\begin{table}[ht]
\centering
\scriptsize
\renewcommand{\arraystretch}{1.18}
\begin{tabular}{>{\raggedright\arraybackslash}p{0.20\textwidth}
                >{\raggedright\arraybackslash}p{0.33\textwidth}
                >{\raggedright\arraybackslash}p{0.38\textwidth}}
\toprule
Present statement & Prior labelled-path results & Exact comparison\\
\midrule
Scope-indexed \(\Ren/\Sub\) algebra
& The 2016 and 2017 papers develop labelled equality reasons and their
  identity-type reading; the 2018 paper makes path operations explicit
  \cite{deQueirozEtAl2016,RamosEtAl2017,RamosEtAl2018}.
& No absence claim is made about substitution results in those works.
  This paper isolates and proves the specific same-scope de Bruijn
  \(\Ren(n,m)\), \(\Sub(n,m)\), double-lift, and instantiation interface used
  by its theorem package.\\

Ten top-level primitive-redex cases
& The prior calculus studies equality rewrites and computational-path
  reductions; the monograph gives a comprehensive labelled calculus
  \cite{RamosEtAl2026}.
& Theorem~\ref{thm:redex-preservation} is narrower: ten raw constructors,
  no contextual closure, no multi-step theorem, and only untyped
  \(B\defeq A\) coherence between displayed types.\\

Source identity rewrite soundness
& Groupoid and rewrite laws for label-sensitive computational paths are
  central in the earlier papers and book
  \cite{RamosEtAl2017,RamosEtAl2018,deVerasEtAl2025,RamosEtAl2026}.
& Theorem~\ref{thm:identity-soundness} is a separation result for the
  present source grammar.  Its target is only the eight-rule core plus
  congruence on unlabelled quotient traces.\\

Endpoint exactness
& Earlier work relates labelled paths and equality/identity judgments
  \cite{deQueirozEtAl2016,RamosEtAl2017}.
& Proposition~\ref{prop:intro-endpoint} is not a strengthening of those
  results.  It is tautological exactness for the quotient by \(\defeq\).\\

Metadata-fiber criterion and trace obstruction
& Earlier identity-type discussions concern the labelled computational-path
  programme; the book develops the broader calculus
  \cite{RamosEtAl2017,RamosEtAl2026}.
& The standard background total-space contraction is refined by
  Theorem~\ref{thm:metadata-fiber-criterion}:
  an unrestricted based eliminator satisfying propositional beta for \(R_M\)
  exists exactly when
  \(M(a,\mathsf{refl})\) is contractible.
  Theorem~\ref{thm:no-trace-j} applies the criterion to visible lists; it
  does not preclude an unrestricted based eliminator with propositional beta
  for a quotient or hidden-label identity type.\\

Information carried by a path
& A labelled path retains a computational reason and rewrites that reason
  \cite{RamosEtAl2018,RamosEtAl2026}.
& A trace atom here stores only
  \((\den t,\den u,\bar h)\); the source derivation \(h\) is erased.
  Evaluated programs decorate already quotient-collapsed endpoints.\\
\bottomrule
\end{tabular}
\caption{Statement-by-statement comparison with the labelled
computational-path literature.}
\label{tab:comparison}
\end{table}

\section{Related work}
\label{sec:related}

\subsection{Syntax and categorical semantics of dependent type theory}

The scope representation follows de Bruijn's nameless notation
\cite{deBruijn1972}; the judgmental organization follows the intensional
tradition of Martin-L\"of~\cite{MartinLof1984}.  Contextual categories and generalized
algebraic theories provide a categorical account of contexts and
substitution~\cite{Cartmell1986}.  Dybjer's internal type theory
\cite{Dybjer1996} and Hofmann's syntax--semantics account
\cite{Hofmann1997} are particularly relevant to the separation between raw
substitution and semantic substitution.  Uemura's general framework
\cite{Uemura2023} gives a much broader organization of theories and their
models.  We do not establish the initiality or universal-model properties
studied in that line of work; \(\Q_n\) is only an explicitly constructed
syntactic quotient.

\subsection{Identity types and homotopical semantics}

The Hofmann--Streicher groupoid interpretation
\cite{HofmannStreicher1998} demonstrated that intensional identity types need
not collapse to proof-irrelevant equality.  Homotopy-theoretic models
\cite{AwodeyWarren2009} and the identity-type weak factorization system
\cite{GambinoGarner2008} place identity elimination in categorical and
homotopical structure.  Lumsdaine~\cite{Lumsdaine2010} and van den Berg and
Garner~\cite{vanDenBergGarner2011} establish weak higher-groupoid structure
generated by intensional identity types.  The HoTT Book
\cite{HoTTBook2013} develops this perspective together with univalence and
higher inductive types.

The familiar path-induction proof that the based total identity space
\(\sum_{b:A}(a=b)\) is contractible \cite{HoTTBook2013} is the starting point
for the metadata-fiber criterion of \cite{ObservableMetadata2026}, which
identifies the total space of an equality-plus-metadata refinement with its
reflexivity metadata fiber.  Section~\ref{sec:j-boundary} applies that
criterion to quotient-trace records.

Our target is different in a decisive respect: it is a record consisting of
proof-irrelevant equality and an observable trace list.  The no-go theorem
in Section~\ref{sec:j-boundary} explains why the higher-groupoid results for
ordinary identity types cannot simply be transferred to this record.  The
sound groupoid rewrites of Section~\ref{sec:identity-programs} concern
evaluations of source identity programs, not a derivation of an
identity-type eliminator.

\subsection{Strict and two-level equality}

Two-level systems distinguish a strict equality suitable for coherence from
a fibrant or homotopical equality.  Capriotti's models
\cite{Capriotti2016}, the strict-equality extension of Altenkirch, Capriotti,
and Kraus~\cite{AltenkirchCapriottiKraus2016}, and the systematic treatment
by Annenkov, Capriotti, Kraus, and Sattler~\cite{AnnenkovEtAl2023} provide
the relevant comparison.  Our layered table is not itself a two-level type
theory: ambient equality, source definitional equality, source identity
programs, and target trace rewrites are metatheoretic layers, not two
internal equality formers with a fibrancy discipline.  Nevertheless, the
two-level literature reinforces the methodological lesson that equality
notions and their elimination rules must be kept explicit.

\subsection{Computational paths}

The 2016 and 2017 articles
\cite{deQueirozEtAl2016,RamosEtAl2017} connect labelled equality reasons
with identity types; the 2018 treatment makes computational paths explicit
\cite{RamosEtAl2018}; and later work develops weak groupoid structure
\cite{deVerasEtAl2025}.  The monograph
\cite{RamosEtAl2026} is the comprehensive background source for the
label-sensitive calculus, its history, rewriting operations, and
applications.  The present target deliberately erases those labels at
quotient soundness.  The structural results and this change of target are
summarized in Table~\ref{tab:comparison}; no claim of absolute priority is
intended.
\section{Reproducibility}
\label{sec:reproducibility}

The mathematical proofs were checked against a machine-verifiable
development in Lean~4~\cite{deMouraUllrich2021,ScopedTraceSoftware2026}.
The artifact covers the scoped syntax, structural laws, raw judgments,
primitive-redex preservation, contextual and multi-step reduction, the failure
and inadmissibility theorems, the conversion-closed extension, quotient
substitution, erasure with its factorization and essential image, source identity rewrite soundness, endpoint exactness, the
independent fixed-universe algebra, the abstract total-space equivalence, the
metadata-fiber classification, and the concrete trace-metadata obstruction.

Every displayed theorem was checked with Lean's axiom report.  The
machine-checked results rely only on propositional extensionality and quotient
soundness, and the negative theorems of Section~\ref{sec:reduction} use
propositional extensionality alone; no form of choice is used, so the results
are constructive relative to those principles.  Two points deserve emphasis.
The conversion hypothesis of Definition~\ref{def:conversion-rules} is formalized
as a record of assumed rules that every dependent theorem takes as an explicit
argument, not as an axiom---the development declares none.  And because
Theorem~\ref{thm:not-admissible} proves that record uninhabited, a reader
auditing the artifact should read the conditional subject-reduction theorem as an
analysis of the induction rather than as a usable lemma.

The theorem-by-theorem correspondence with the source modules, and the build
instructions, are collected in Appendix~\ref{app:artifact}, so that no
formalization detail interrupts the mathematics.  The software is evidence for
the listed inductions; it does not supply the missing context-formation,
confluence, or external interpretation theorems.

\section{Conclusion}
\label{sec:conclusion}

An audit is judged by whether a reader can act on it, so we close by saying
what has been settled and what a user of the artifact should conclude.

On the positive side, the calculus genuinely supports a scoped syntax and its
binder algebra, typed substitution, preservation for ten top-level primitive
redexes, a contextual and multi-step reduction theory, compositional
substitution on a syntactic quotient, and sound evaluation of source identity
rewrites into an explicitly presented unlabelled trace relation.  These are
usable, and none of them requires calling the quotient a model of MLTT.

Two further results are measurements rather than constructions, and they are the
audit's findings.

The first concerns reduction, and it is negative.  Contextual subject reduction
is not merely unproved here but false, with a witness small enough to check by
eye; neither of the two rules that would repair it is admissible; and the
hypothesis record packaging them is uninhabited, so the conditional theorem
derived from it is vacuous as a statement about \(\vdash\).  What that
conditional proof does establish is an analysis of the induction: repairs of
exactly two shapes, enumerated in Table~\ref{tab:sr-cases}, close every case.
The constructive response is the conversion-closed extension \(\vdashC\), for
which we prove no subject-reduction theorem, because doing so needs generation
lemmas and hence confluence.  The distance from this calculus to subject
reduction is therefore a different judgment, not two hypotheses.  For a reader,
the operational consequence is concrete: the raw judgment must not be relied on
to be preserved by reduction anywhere in the development.

The second concerns labels.  Calling the quotient traces \emph{unlabelled} is
otherwise a statement about intent, so Section~\ref{sec:erasure} makes it a
theorem.  Evaluation factors through a label-free syntax whose atoms carry only
proof-irrelevant equalities of classes; erasure preserves program size, so the
arrangement of equality steps survives and only labels are discarded; and at
atoms erasure is inverse to the recovery map supplied by quotient exactness.
Because the label-free grammar is strictly more permissive than erasure needs,
we characterize the essential image rather than claiming the two languages
agree.  Together these pin down what the target retains and what it cannot
express, and they make the contrast with label-sensitive computational paths
precise: an analogue of Corollary~\ref{cor:no-derivation-observation} fails
there by design, and that failure is the source of that theory's content.
Applying the metadata-fiber criterion of \cite{ObservableMetadata2026} to our
own records completes the picture quantitatively: the quotient trace fiber is a
list of quotient classes, hence never contractible, so trace records support no
unrestricted based eliminator, while equality-component motives remain available
through projection.  These objects are observable list metadata decorating an
equality, not carriers of reasons for it.  Endpoint exactness, by contrast, is
immediate from quotient soundness and exactness; its role is bookkeeping, not
semantic completeness.

The open directions follow from the above, and we state them as the work the
audit has scoped rather than as speculation.  Confluence for \(\defeq\) would
unlock generation lemmas and hence an unconditional subject-reduction theorem
for \(\vdashC\); that is the single most valuable next step, and this paper
establishes both that it is needed and that nothing weaker will do.  We prove no
normalization or standardization theorem, the syntactic quotient is built from
\(\defeq\) directly rather than from a normal-form analysis, and a genuine
raw-to-CwF interpretation is not attempted.  Each such extension must state
which equality layer and which elimination principle it uses.

\appendix

\section{The companion Lean artifact}
\label{app:artifact}

The correspondence between the theorems of the article and the development is as
follows.  Contextual and multi-step reduction, together with their soundness
and substitution stability
(Definition~\ref{def:contextual-reduction},
Proposition~\ref{prop:reduction-basic}), are
\path{Reduction}, \path{ReductionMany}, \path{Reduction.toDefEq},
\path{ReductionMany.toDefEq}, \path{Reduction.substitution}, and
\path{ReductionMany.substitution}.  The conversion hypothesis of
Definition~\ref{def:conversion-rules} is the record \path{ConversionRules};
Lemma~\ref{lem:ctx-congruence} is \path{ctxDefEq_extend},
\path{natStepCtx_defEq}, \path{eqMotiveCtx_defEq},
\path{instantiate_congr_body} and its two-variable analogue; and
Theorem~\ref{thm:contextual-subject-reduction} together with
Corollary~\ref{cor:multistep-subject-reduction} are
\path{contextual_subject_reduction} and \path{multistep_subject_reduction}.

The erasure results of Section~\ref{sec:erasure} are
\path{QuotProgram} and \path{erase} for
Definitions~\ref{def:quot-programs} and~\ref{def:erasure};
\path{eval_factors_through_erase} for
Theorem~\ref{thm:erasure-factorization};
\path{defEqPath_derivation_irrelevant} and
\path{atom_eval_derivation_irrelevant} for
Corollary~\ref{cor:no-derivation-observation};
\path{erase_preserves_size} for Proposition~\ref{prop:erasure-size};
\path{source_atom_iff_quotient_atom} and \path{erase_atom_roundTrip} for
Proposition~\ref{prop:erasure-roundtrip}; and
\path{erased_identity_rweq_sound} for
Corollary~\ref{cor:erased-soundness}.  The computed trace fiber and the
resulting obstruction are \path{quotientTraceEquivClassList} and
\path{quotient_trace_fiber_not_contractible}.

The negative results of Section~\ref{sec:reduction} are
\path{failSource_typed}, \path{failSource_reduces} and
\path{failTarget_untyped} for Theorem~\ref{thm:sr-fails}, assembled as
\path{raw_contextual_subject_reduction_fails};
\path{convert_not_admissible} and \path{convertCtx_not_admissible} for
parts~(i) and~(ii) of Theorem~\ref{thm:not-admissible}, with
\path{failFamily_typed} witnessing that the family in~(ii) is a type over the
unreduced domain; and \path{conversionRules_uninhabited} for part~(iii).  The
conversion-closed extension of Definition~\ref{def:vdashC} is \path{HasTypeC},
with \path{hasTypeC_convert}, \path{hasTypeC_convertCtx},
\path{HasTypeC.ofRaw} and \path{failTarget_typedC} for
Proposition~\ref{prop:vdashC}.  The image characterization of
Theorem~\ref{thm:erasure-image} is \path{FramedAt}, \path{erase_framedAt},
\path{framedAt_surjective} and \path{erase_image_eq_framedAt}, with
\path{erase_congr_frame_induced} and \path{erase_atom_is_sound} for the two
``invents nothing'' statements.

The axiom profile of the development, and the reason
\path{contextual_subject_reduction} must be read as an analysis of the induction
rather than as a usable lemma, are stated in Section~\ref{sec:reproducibility}
and not repeated here.

\medskip
\noindent
\fbox{\parbox{0.94\linewidth}{
\textbf{Artifact.}
Repository: \url{https://github.com/Arthur742Ramos/ComputationalPathsLean}.
Toolchain: Lean~4 (\texttt{leanprover/lean4:v4.24.0}), pinned in
\path{lean-toolchain}.  Dependencies: \path{mathlib4} at tag \texttt{v4.24.0},
pinned with the full transitive lock in \path{lake-manifest.json}.
Build: \texttt{lake build}.  Sources for this article:
\path{ComputationalPaths/Path/TypeTheory/Raw*.lean}.
Invariant check (no \texttt{sorry}, \texttt{admit}, or custom \texttt{axiom}):
\texttt{python3 scripts/check\_invariants.py}.
Axiom audit: \texttt{\#print axioms} on any theorem named in this appendix.
A Zenodo archival DOI will be minted for the version of record; until then the
tagged commit in the repository is the citable snapshot.}}

\medskip
The software is evidence for the listed inductions; it does not supply the
missing context-formation, confluence, or external interpretation theorems.

\small

\section{Raw grammar and judgment summary}
\label{app:grammar}

\subsection{Scope-indexed constructors}

The complete scope information is as follows.  In each row, unmarked
arguments live in \(\Expr_n\).
\[
\begin{array}{rcl}
\Var &:& \Fin(n)\to\Expr_n,\\
\Sort_\ell&:&\Expr_n,\\
\PiTy,\SigmaTy&:&\Expr_n\times\Expr_{n+1}\to\Expr_n,\\
\Lam&:&\Expr_{n+1}\to\Expr_n,\\
\App,\Pair&:&\Expr_n\times\Expr_n\to\Expr_n,\\
\Fst,\Snd&:&\Expr_n\to\Expr_n,\\
\NatTy_\ell,\Zero_\ell&:&\Expr_n,\\
\Succ_\ell&:&\Expr_n\to\Expr_n,\\
\NatElim_\ell&:&
  \Expr_{n+1}\times\Expr_n\times\Expr_{n+2}\times\Expr_n
  \to\Expr_n,\\
\CodeNat_\ell&:&\Expr_n,\\
\CodePi,\CodeSigma&:&\Expr_n\times\Expr_{n+1}\to\Expr_n,\\
\CodeId&:&\Expr_n^3\to\Expr_n,\\
\El&:&\Expr_n\to\Expr_n,\\
\IdTy&:&\Expr_n^3\to\Expr_n,\\
\Refl&:&\Expr_n\to\Expr_n,\\
\EqJ&:&\Expr_{n+2}\times\Expr_n^3\to\Expr_n.
\end{array}
\tag{A.1}\label{eq:full-grammar}
\]

\subsection{Renaming and substitution clauses under binders}

The nontrivial recursion clauses are
\[
\begin{aligned}
\PiTy(A,B)[\rho]_{\mathrm r}
 &=\PiTy(A[\rho]_{\mathrm r},B[\rho^\lift]_{\mathrm r}),\\
\Lam(t)[\rho]_{\mathrm r}
 &=\Lam(t[\rho^\lift]_{\mathrm r}),\\
\NatElim_\ell(M,z,s,k)[\rho]_{\mathrm r}
 &=\NatElim_\ell(M[\rho^\lift]_{\mathrm r},
    z[\rho]_{\mathrm r},s[\rho^{\lift\lift}]_{\mathrm r},
    k[\rho]_{\mathrm r}),\\
\EqJ(M,r,b,p)[\rho]_{\mathrm r}
 &=\EqJ(M[\rho^{\lift\lift}]_{\mathrm r},
    r[\rho]_{\mathrm r},b[\rho]_{\mathrm r},p[\rho]_{\mathrm r}),
\end{aligned}
\]
with identical lift placement for substitution:
\[
\begin{aligned}
\PiTy(A,B)[\sigma]
 &=\PiTy(A[\sigma],B[\sigma^\lift]),\\
\Lam(t)[\sigma]&=\Lam(t[\sigma^\lift]),\\
\NatElim_\ell(M,z,s,k)[\sigma]
 &=\NatElim_\ell(M[\sigma^\lift],z[\sigma],
     s[\sigma^{\lift\lift}],k[\sigma]),\\
\EqJ(M,r,b,p)[\sigma]
 &=\EqJ(M[\sigma^{\lift\lift}],r[\sigma],
        b[\sigma],p[\sigma]).
\end{aligned}
\tag{A.2}\label{eq:binder-clauses}
\]
Product and sum codes use one lift in their second argument.

\subsection{Typing rules omitted from the main display}

The complete sum fragment is
\[
\frac{\Gamma\vdash A:\Sort_\ell\quad
      \Ext{\Gamma}{A}\vdash B:\Sort_\ell}
     {\Gamma\vdash\SigmaTy(A,B):\Sort_\ell},
\qquad
\frac{\Gamma\vdash a:A\quad\Gamma\vdash b:B\{a\}}
     {\Gamma\vdash\Pair(a,b):\SigmaTy(A,B)},
\]
\[
\frac{\Gamma\vdash p:\SigmaTy(A,B)}
     {\Gamma\vdash\Fst(p):A},
\qquad
\frac{\Gamma\vdash p:\SigmaTy(A,B)}
     {\Gamma\vdash\Snd(p):B\{\Fst(p)\}}.
\tag{A.3}\label{eq:sigma-rules}
\]

For natural elimination, define the successor-branch context
\[
 \Gamma_{\mathrm{succ}}
   =\Ext{\Ext{\Gamma}{\NatTy_\ell}}{M}.
\]
In its scope, let
\[
 T_{\mathrm{succ}}(M)
   =M[
      0\mapsto\Succ_\ell(\Var(1)),\
      i+1\mapsto\Var(i+2)
     ].
\tag{A.4}\label{eq:nat-step-target}
\]
If
\[
\Ext{\Gamma}{\NatTy_\ell}\vdash M:\Sort_j,\quad
\Gamma\vdash z:M\{\Zero_\ell\},\quad
\Gamma_{\mathrm{succ}}\vdash s:T_{\mathrm{succ}}(M),\quad
\Gamma\vdash k:\NatTy_\ell,
\]
then
\[
\Gamma\vdash\NatElim_\ell(M,z,s,k):M\{k\}.
\tag{A.5}\label{eq:nat-elim-full}
\]
Two-variable instantiation of~\eqref{eq:nat-step-target} by \(k\) and an
induction hypothesis is \(M\{\Succ_\ell(k)\}\), the key equation in the
successor primitive-redex preservation case.

Code formation and decoding are
\[
\frac{}{\Gamma\vdash\CodeNat_\ell:\Sort_\ell},
\quad
\frac{\Gamma\vdash A:\Sort_\ell\quad
      \Ext{\Gamma}{\El(A)}\vdash B:\Sort_\ell}
     {\Gamma\vdash\CodePi(A,B):\Sort_\ell},
\]
\[
\frac{\Gamma\vdash A:\Sort_\ell\quad
      \Ext{\Gamma}{\El(A)}\vdash B:\Sort_\ell}
     {\Gamma\vdash\CodeSigma(A,B):\Sort_\ell},
\]
\[
\frac{\Gamma\vdash A:\Sort_\ell\quad
      \Gamma\vdash a:\El(A)\quad\Gamma\vdash b:\El(A)}
     {\Gamma\vdash\CodeId(A,a,b):\Sort_\ell},
\qquad
\frac{\Gamma\vdash c:\Sort_\ell}
     {\Gamma\vdash\El(c):\Sort_\ell}.
\tag{A.6}\label{eq:code-rules}
\]
The identity rules and raw two-binder eliminator are those in
\eqref{eq:eqj-rule}.

\subsection{Definitional-equality compatible closure}

Besides equivalence and the beta inclusions, \(\defeq\) is congruent in the
following positions:
\[
\begin{array}{ll}
\PiTy,\SigmaTy:&\text{domain and codomain},\\
\Lam:&\text{body},\\
\App:&\text{function and argument},\\
\Pair:&\text{first and second component},\\
\Fst,\Snd,\Succ_\ell,\El,\Refl:&\text{their unique argument},\\
\NatElim_\ell:&\text{motive, zero case, successor case, scrutinee},\\
\CodePi,\CodeSigma:&\text{domain and codomain},\\
\CodeId,\IdTy:&\text{type, left endpoint, right endpoint},\\
\EqJ:&\text{motive, reflexivity case, endpoint, proof}.
\end{array}
\tag{A.7}\label{eq:defeq-congruence}
\]
Binder positions live in their larger scopes.  Substitution closure applies
\(\sigma^\lift\) or \(\sigma^{\lift\lift}\) in precisely those positions.

\section{Source identity rewrite rules}
\label{app:rewrites}

\subsection{Source one-step judgment}

The judgment \(p\srw q\) is indexed by common source endpoints:
\[
 p,q:\IProg_n(t,u)\quad\Longrightarrow\quad
 (p\srw q)\ \mathsf{judgment}.
\]
For \(p:\IProg_n(t,u)\), \(q:\IProg_n(u,v)\), and
\(r:\IProg_n(v,w)\), its eight base constructors are
\[
\frac{}{(\mathbf1_t)^{-1}\srw\mathbf1_t}
\quad
\frac{}{(p^{-1})^{-1}\srw p}
\quad
\frac{}{\mathbf1_t\cdot p\srw p}
\quad
\frac{}{p\cdot\mathbf1_u\srw p},
\tag{B.1}\label{eq:app-source-units}
\]
\[
\frac{}{p\cdot p^{-1}\srw\mathbf1_t}
\quad
\frac{}{p^{-1}\cdot p\srw\mathbf1_u}
\quad
\frac{}{(p\cdot q)^{-1}\srw q^{-1}\cdot p^{-1}}
\quad
\frac{}{(p\cdot q)\cdot r\srw p\cdot(q\cdot r)}.
\tag{B.2}\label{eq:app-source-groupoid}
\]
The compatible constructors are
\[
\frac{p\srw p'}{p^{-1}\srw(p')^{-1}},
\qquad
\frac{p\srw p'}{p\cdot q\srw p'\cdot q},
\qquad
\frac{q\srw q'}{p\cdot q\srw p\cdot q'},
\tag{B.3}\label{eq:app-source-composition-congruence}
\]
\[
\frac{p\srw p'}{C[p]\srw C[p']},
\qquad
\frac{}{C[\mathbf1_t]\srw\mathbf1_{C[t]}}
\quad(C\in\Frame_n).
\tag{B.4}\label{eq:app-source-frame-congruence}
\]
No other source one-step constructor is implicit.

\subsection{Source rewrite equivalence}

At fixed \(t,u\), the judgment \(p\srweq q\) is inductively generated by
\[
\frac{}{p\srweq p},
\qquad
\frac{p\srw q}{p\srweq q},
\qquad
\frac{p\srweq q}{q\srweq p},
\qquad
\frac{p\srweq q\quad q\srweq r}{p\srweq r}.
\tag{B.5}\label{eq:app-source-equivalence}
\]
Thus \(\srweq\) is a derivation-carrying symmetric, reflexive, transitive
closure, not an equality between derivations.

\subsection{Target trace rules}

For \(S:\Tr_X(a,b)\), \(T:\Tr_X(b,c)\), and
\(U:\Tr_X(c,d)\), target one-step rewriting has exactly
\[
\frac{}{(\mathbf r_a)^{-1}\TrStep\mathbf r_a}
\quad
\frac{}{(S^{-1})^{-1}\TrStep S}
\quad
\frac{}{\mathbf r_a\cdot S\TrStep S}
\quad
\frac{}{S\cdot\mathbf r_b\TrStep S},
\tag{B.6}\label{eq:app-target-units}
\]
\[
\frac{}{S\cdot S^{-1}\TrStep\mathbf r_a}
\quad
\frac{}{S^{-1}\cdot S\TrStep\mathbf r_b}.
\]
\[
\frac{}{(S\cdot T)^{-1}\TrStep T^{-1}\cdot S^{-1}}
\quad
\frac{}{(S\cdot T)\cdot U\TrStep S\cdot(T\cdot U)}.
\tag{B.7}\label{eq:app-target-groupoid}
\]
Target equivalence \(S\TrEq T\) is generated by
\[
\frac{}{S\TrEq S},
\quad
\frac{S\TrStep T}{S\TrEq T},
\quad
\frac{S\TrEq T}{T\TrEq S},
\quad
\frac{S\TrEq T\quad T\TrEq U}{S\TrEq U},
\tag{B.8}\label{eq:app-target-equivalence}
\]
together with all four congruence rules
\[
\frac{S\TrEq T}{S^{-1}\TrEq T^{-1}},
\qquad
\frac{S\TrEq T}{S\cdot U\TrEq T\cdot U},
\qquad
\frac{T\TrEq U}{S\cdot T\TrEq S\cdot U},
\qquad
\frac{S\TrEq T}{f_\ast S\TrEq f_\ast T}.
\tag{B.9}\label{eq:app-target-congruence}
\]
The map operation \(f_\ast\) is
Definition~\ref{def:trace-record} and equation~\eqref{eq:trace-map};
in particular \(f_\ast\mathbf r_a=\mathbf r_{fa}\) by definition.

\subsection{Rule-by-rule soundness map}

Evaluation sends the four rules in~\eqref{eq:app-source-units} to
\eqref{eq:app-target-units}, the four rules
in~\eqref{eq:app-source-groupoid} to
\eqref{eq:app-target-groupoid}, and the first three rules
in~\eqref{eq:app-source-composition-congruence} to the first three congruence
rules in~\eqref{eq:app-target-congruence}.  Frame congruence uses the last
target congruence rule with \(f=\overline C\); frame-unit soundness is target
reflexivity after the definitional equation
\(\overline C_\ast\mathbf r_{\den t}=\mathbf r_{\den{C[t]}}\).
Finally, the four constructors in~\eqref{eq:app-source-equivalence} map to
the four constructors in~\eqref{eq:app-target-equivalence}.  This is the
complete induction for Theorem~\ref{thm:identity-soundness}.

\begingroup
\small
\bibliographystyle{amsplain}
\bibliography{refs}
\endgroup

\end{document}
